% figs.tex —— 由 ch12/make_figs.py 產生，請勿手動修改。
% 重新產生：python ch12/make_figs.py（第 12 章，動畫 W1201–W1206）

\newcommand{\TwelveFigShell}{%
\begin{tikzpicture}[font=\scriptsize,>=Stealth,baseline=(current bounding box.center)]
\fill[cWalk!22,even odd rule] (0,0) circle (1.450) (0,0) circle (1.096);
\draw[cBar,thick] (0,0) circle (1.450);
\draw[cBar,densely dashed] (0,0) circle (1.096);
\fill[cWalk] (-1.40,-0.20) circle (0.7pt);
\fill[cWalk] (-1.40,0.00) circle (0.7pt);
\fill[cWalk] (-1.40,0.20) circle (0.7pt);
\fill[cWalk] (-1.20,-0.80) circle (0.7pt);
\fill[cWalk] (-1.20,-0.60) circle (0.7pt);
\fill[cWalk] (-1.20,-0.40) circle (0.7pt);
\fill[cWalk] (-1.20,-0.20) circle (0.7pt);
\fill[cWalk] (-1.20,0.00) circle (0.7pt);
\fill[cWalk] (-1.20,0.20) circle (0.7pt);
\fill[cWalk] (-1.20,0.40) circle (0.7pt);
\fill[cWalk] (-1.20,0.60) circle (0.7pt);
\fill[cWalk] (-1.20,0.80) circle (0.7pt);
\fill[cWalk] (-1.00,-1.00) circle (0.7pt);
\fill[cWalk] (-1.00,-0.80) circle (0.7pt);
\fill[cWalk] (-1.00,-0.60) circle (0.7pt);
\fill[cBar!60] (-1.00,-0.40) circle (0.7pt);
\fill[cBar!60] (-1.00,-0.20) circle (0.7pt);
\fill[cBar!60] (-1.00,0.00) circle (0.7pt);
\fill[cBar!60] (-1.00,0.20) circle (0.7pt);
\fill[cBar!60] (-1.00,0.40) circle (0.7pt);
\fill[cWalk] (-1.00,0.60) circle (0.7pt);
\fill[cWalk] (-1.00,0.80) circle (0.7pt);
\fill[cWalk] (-1.00,1.00) circle (0.7pt);
\fill[cWalk] (-0.80,-1.20) circle (0.7pt);
\fill[cWalk] (-0.80,-1.00) circle (0.7pt);
\fill[cWalk] (-0.80,-0.80) circle (0.7pt);
\fill[cBar!60] (-0.80,-0.60) circle (0.7pt);
\fill[cBar!60] (-0.80,-0.40) circle (0.7pt);
\fill[cBar!60] (-0.80,-0.20) circle (0.7pt);
\fill[cBar!60] (-0.80,0.00) circle (0.7pt);
\fill[cBar!60] (-0.80,0.20) circle (0.7pt);
\fill[cBar!60] (-0.80,0.40) circle (0.7pt);
\fill[cBar!60] (-0.80,0.60) circle (0.7pt);
\fill[cWalk] (-0.80,0.80) circle (0.7pt);
\fill[cWalk] (-0.80,1.00) circle (0.7pt);
\fill[cWalk] (-0.80,1.20) circle (0.7pt);
\fill[cWalk] (-0.60,-1.20) circle (0.7pt);
\fill[cWalk] (-0.60,-1.00) circle (0.7pt);
\fill[cBar!60] (-0.60,-0.80) circle (0.7pt);
\fill[cBar!60] (-0.60,-0.60) circle (0.7pt);
\fill[cBar!60] (-0.60,-0.40) circle (0.7pt);
\fill[cBar!60] (-0.60,-0.20) circle (0.7pt);
\fill[cBar!60] (-0.60,0.00) circle (0.7pt);
\fill[cBar!60] (-0.60,0.20) circle (0.7pt);
\fill[cBar!60] (-0.60,0.40) circle (0.7pt);
\fill[cBar!60] (-0.60,0.60) circle (0.7pt);
\fill[cBar!60] (-0.60,0.80) circle (0.7pt);
\fill[cWalk] (-0.60,1.00) circle (0.7pt);
\fill[cWalk] (-0.60,1.20) circle (0.7pt);
\fill[cWalk] (-0.40,-1.20) circle (0.7pt);
\fill[cBar!60] (-0.40,-1.00) circle (0.7pt);
\fill[cBar!60] (-0.40,-0.80) circle (0.7pt);
\fill[cBar!60] (-0.40,-0.60) circle (0.7pt);
\fill[cBar!60] (-0.40,-0.40) circle (0.7pt);
\fill[cBar!60] (-0.40,-0.20) circle (0.7pt);
\fill[cBar!60] (-0.40,0.00) circle (0.7pt);
\fill[cBar!60] (-0.40,0.20) circle (0.7pt);
\fill[cBar!60] (-0.40,0.40) circle (0.7pt);
\fill[cBar!60] (-0.40,0.60) circle (0.7pt);
\fill[cBar!60] (-0.40,0.80) circle (0.7pt);
\fill[cBar!60] (-0.40,1.00) circle (0.7pt);
\fill[cWalk] (-0.40,1.20) circle (0.7pt);
\fill[cWalk] (-0.20,-1.40) circle (0.7pt);
\fill[cWalk] (-0.20,-1.20) circle (0.7pt);
\fill[cBar!60] (-0.20,-1.00) circle (0.7pt);
\fill[cBar!60] (-0.20,-0.80) circle (0.7pt);
\fill[cBar!60] (-0.20,-0.60) circle (0.7pt);
\fill[cBar!60] (-0.20,-0.40) circle (0.7pt);
\fill[cBar!60] (-0.20,-0.20) circle (0.7pt);
\fill[cBar!60] (-0.20,0.00) circle (0.7pt);
\fill[cBar!60] (-0.20,0.20) circle (0.7pt);
\fill[cBar!60] (-0.20,0.40) circle (0.7pt);
\fill[cBar!60] (-0.20,0.60) circle (0.7pt);
\fill[cBar!60] (-0.20,0.80) circle (0.7pt);
\fill[cBar!60] (-0.20,1.00) circle (0.7pt);
\fill[cWalk] (-0.20,1.20) circle (0.7pt);
\fill[cWalk] (-0.20,1.40) circle (0.7pt);
\fill[cWalk] (0.00,-1.40) circle (0.7pt);
\fill[cWalk] (0.00,-1.20) circle (0.7pt);
\fill[cBar!60] (0.00,-1.00) circle (0.7pt);
\fill[cBar!60] (0.00,-0.80) circle (0.7pt);
\fill[cBar!60] (0.00,-0.60) circle (0.7pt);
\fill[cBar!60] (0.00,-0.40) circle (0.7pt);
\fill[cBar!60] (0.00,-0.20) circle (0.7pt);
\fill[cBar!60] (0.00,0.00) circle (0.7pt);
\fill[cBar!60] (0.00,0.20) circle (0.7pt);
\fill[cBar!60] (0.00,0.40) circle (0.7pt);
\fill[cBar!60] (0.00,0.60) circle (0.7pt);
\fill[cBar!60] (0.00,0.80) circle (0.7pt);
\fill[cBar!60] (0.00,1.00) circle (0.7pt);
\fill[cWalk] (0.00,1.20) circle (0.7pt);
\fill[cWalk] (0.00,1.40) circle (0.7pt);
\fill[cWalk] (0.20,-1.40) circle (0.7pt);
\fill[cWalk] (0.20,-1.20) circle (0.7pt);
\fill[cBar!60] (0.20,-1.00) circle (0.7pt);
\fill[cBar!60] (0.20,-0.80) circle (0.7pt);
\fill[cBar!60] (0.20,-0.60) circle (0.7pt);
\fill[cBar!60] (0.20,-0.40) circle (0.7pt);
\fill[cBar!60] (0.20,-0.20) circle (0.7pt);
\fill[cBar!60] (0.20,0.00) circle (0.7pt);
\fill[cBar!60] (0.20,0.20) circle (0.7pt);
\fill[cBar!60] (0.20,0.40) circle (0.7pt);
\fill[cBar!60] (0.20,0.60) circle (0.7pt);
\fill[cBar!60] (0.20,0.80) circle (0.7pt);
\fill[cBar!60] (0.20,1.00) circle (0.7pt);
\fill[cWalk] (0.20,1.20) circle (0.7pt);
\fill[cWalk] (0.20,1.40) circle (0.7pt);
\fill[cWalk] (0.40,-1.20) circle (0.7pt);
\fill[cBar!60] (0.40,-1.00) circle (0.7pt);
\fill[cBar!60] (0.40,-0.80) circle (0.7pt);
\fill[cBar!60] (0.40,-0.60) circle (0.7pt);
\fill[cBar!60] (0.40,-0.40) circle (0.7pt);
\fill[cBar!60] (0.40,-0.20) circle (0.7pt);
\fill[cBar!60] (0.40,0.00) circle (0.7pt);
\fill[cBar!60] (0.40,0.20) circle (0.7pt);
\fill[cBar!60] (0.40,0.40) circle (0.7pt);
\fill[cBar!60] (0.40,0.60) circle (0.7pt);
\fill[cBar!60] (0.40,0.80) circle (0.7pt);
\fill[cBar!60] (0.40,1.00) circle (0.7pt);
\fill[cWalk] (0.40,1.20) circle (0.7pt);
\fill[cWalk] (0.60,-1.20) circle (0.7pt);
\fill[cWalk] (0.60,-1.00) circle (0.7pt);
\fill[cBar!60] (0.60,-0.80) circle (0.7pt);
\fill[cBar!60] (0.60,-0.60) circle (0.7pt);
\fill[cBar!60] (0.60,-0.40) circle (0.7pt);
\fill[cBar!60] (0.60,-0.20) circle (0.7pt);
\fill[cBar!60] (0.60,0.00) circle (0.7pt);
\fill[cBar!60] (0.60,0.20) circle (0.7pt);
\fill[cBar!60] (0.60,0.40) circle (0.7pt);
\fill[cBar!60] (0.60,0.60) circle (0.7pt);
\fill[cBar!60] (0.60,0.80) circle (0.7pt);
\fill[cWalk] (0.60,1.00) circle (0.7pt);
\fill[cWalk] (0.60,1.20) circle (0.7pt);
\fill[cWalk] (0.80,-1.20) circle (0.7pt);
\fill[cWalk] (0.80,-1.00) circle (0.7pt);
\fill[cWalk] (0.80,-0.80) circle (0.7pt);
\fill[cBar!60] (0.80,-0.60) circle (0.7pt);
\fill[cBar!60] (0.80,-0.40) circle (0.7pt);
\fill[cBar!60] (0.80,-0.20) circle (0.7pt);
\fill[cBar!60] (0.80,0.00) circle (0.7pt);
\fill[cBar!60] (0.80,0.20) circle (0.7pt);
\fill[cBar!60] (0.80,0.40) circle (0.7pt);
\fill[cBar!60] (0.80,0.60) circle (0.7pt);
\fill[cWalk] (0.80,0.80) circle (0.7pt);
\fill[cWalk] (0.80,1.00) circle (0.7pt);
\fill[cWalk] (0.80,1.20) circle (0.7pt);
\fill[cWalk] (1.00,-1.00) circle (0.7pt);
\fill[cWalk] (1.00,-0.80) circle (0.7pt);
\fill[cWalk] (1.00,-0.60) circle (0.7pt);
\fill[cBar!60] (1.00,-0.40) circle (0.7pt);
\fill[cBar!60] (1.00,-0.20) circle (0.7pt);
\fill[cBar!60] (1.00,0.00) circle (0.7pt);
\fill[cBar!60] (1.00,0.20) circle (0.7pt);
\fill[cBar!60] (1.00,0.40) circle (0.7pt);
\fill[cWalk] (1.00,0.60) circle (0.7pt);
\fill[cWalk] (1.00,0.80) circle (0.7pt);
\fill[cWalk] (1.00,1.00) circle (0.7pt);
\fill[cWalk] (1.20,-0.80) circle (0.7pt);
\fill[cWalk] (1.20,-0.60) circle (0.7pt);
\fill[cWalk] (1.20,-0.40) circle (0.7pt);
\fill[cWalk] (1.20,-0.20) circle (0.7pt);
\fill[cWalk] (1.20,0.00) circle (0.7pt);
\fill[cWalk] (1.20,0.20) circle (0.7pt);
\fill[cWalk] (1.20,0.40) circle (0.7pt);
\fill[cWalk] (1.20,0.60) circle (0.7pt);
\fill[cWalk] (1.20,0.80) circle (0.7pt);
\fill[cWalk] (1.40,-0.20) circle (0.7pt);
\fill[cWalk] (1.40,0.00) circle (0.7pt);
\fill[cWalk] (1.40,0.20) circle (0.7pt);
\draw[->] (0,0) -- (1.197,0.819) node[pos=0.55,above left=-1pt] {$\Lambda=e^{-\ell}$};
\node[align=center] at (0,-1.900) {殼層（橙）：$e^{-\ell-\delta\ell}<|\vec q|<e^{-\ell}$};
\end{tikzpicture}}

\newcommand{\TwelveTabShell}{%
\begin{tabular}{@{}ccrrr@{}}
\toprule
$d$ & $\delta\ell$ & $1-e^{-d\,\delta\ell}$ & 數格點 & $d\,\delta\ell$\\
\midrule
3 & 0.1 & 0.2592 & 0.2591 & 0.3000\\
3 & 0.01 & 0.0296 & 0.0296 & 0.0300\\
4 & 0.01 & 0.0392 & -- & 0.0400\\
\bottomrule
\end{tabular}}

\newcommand{\TwelveFigToy}{%
\begin{tikzpicture}[font=\scriptsize,>=Stealth,baseline=(current bounding box.center)]
\draw[->,gray] (0,0) -- (4.050,0) node[right,black] {$s$};
\draw[->,gray] (1.900,0) -- (1.900,2.950) node[above,black] {};
\draw[gray] (0.000,0) -- ++(0,-0.06) node[below,black] {$-2$};
\draw[gray] (0.950,0) -- ++(0,-0.06) node[below,black] {$-1$};
\draw[gray] (1.900,0) -- ++(0,-0.06) node[below,black] {$0$};
\draw[gray] (2.850,0) -- ++(0,-0.06) node[below,black] {$1$};
\draw[gray] (3.800,0) -- ++(0,-0.06) node[below,black] {$2$};
\draw[gray] (1.900,0.900) -- ++(-0.06,0) node[left,black] {$0.25$};
\draw[gray] (1.900,1.800) -- ++(-0.06,0) node[left,black] {$0.5$};
\draw[gray] (1.900,2.700) -- ++(-0.06,0) node[left,black] {$0.75$};
\draw[gray,densely dashed] (0.000,2.016) -- (0.097,1.763) -- (0.195,1.533) -- (0.292,1.326) -- (0.390,1.140) -- (0.487,0.972) -- (0.585,0.822) -- (0.682,0.689) -- (0.779,0.571) -- (0.877,0.466) -- (0.974,0.374) -- (1.072,0.294) -- (1.169,0.226) -- (1.267,0.167) -- (1.364,0.118) -- (1.462,0.078) -- (1.559,0.047) -- (1.656,0.024) -- (1.754,0.009) -- (1.851,0.001) -- (1.949,0.001) -- (2.046,0.009) -- (2.144,0.024) -- (2.241,0.047) -- (2.338,0.078) -- (2.436,0.118) -- (2.533,0.167) -- (2.631,0.226) -- (2.728,0.294) -- (2.826,0.374) -- (2.923,0.466) -- (3.021,0.571) -- (3.118,0.689) -- (3.215,0.822) -- (3.313,0.972) -- (3.410,1.140) -- (3.508,1.326) -- (3.605,1.533) -- (3.703,1.763) -- (3.800,2.016);
\draw[cBar,thick] (0.000,2.622) -- (0.097,2.316) -- (0.195,2.036) -- (0.292,1.780) -- (0.390,1.545) -- (0.487,1.332) -- (0.585,1.138) -- (0.682,0.963) -- (0.779,0.805) -- (0.877,0.663) -- (0.974,0.537) -- (1.072,0.426) -- (1.169,0.329) -- (1.267,0.245) -- (1.364,0.175) -- (1.462,0.116) -- (1.559,0.070) -- (1.656,0.036) -- (1.754,0.013) -- (1.851,0.001) -- (1.949,0.001) -- (2.046,0.013) -- (2.144,0.036) -- (2.241,0.070) -- (2.338,0.116) -- (2.436,0.175) -- (2.533,0.245) -- (2.631,0.329) -- (2.728,0.426) -- (2.826,0.537) -- (2.923,0.663) -- (3.021,0.805) -- (3.118,0.963) -- (3.215,1.138) -- (3.313,1.332) -- (3.410,1.545) -- (3.508,1.780) -- (3.605,2.036) -- (3.703,2.316) -- (3.800,2.622);
\draw[cWalk,thick,densely dotted] (0.000,2.592) -- (0.097,2.294) -- (0.195,2.020) -- (0.292,1.768) -- (0.390,1.537) -- (0.487,1.326) -- (0.585,1.135) -- (0.682,0.961) -- (0.779,0.803) -- (0.877,0.663) -- (0.974,0.537) -- (1.072,0.426) -- (1.169,0.329) -- (1.267,0.245) -- (1.364,0.175) -- (1.462,0.116) -- (1.559,0.070) -- (1.656,0.036) -- (1.754,0.013) -- (1.851,0.001) -- (1.949,0.001) -- (2.046,0.013) -- (2.144,0.036) -- (2.241,0.070) -- (2.338,0.116) -- (2.436,0.175) -- (2.533,0.245) -- (2.631,0.329) -- (2.728,0.426) -- (2.826,0.537) -- (2.923,0.663) -- (3.021,0.803) -- (3.118,0.961) -- (3.215,1.135) -- (3.313,1.326) -- (3.410,1.537) -- (3.508,1.768) -- (3.605,2.020) -- (3.703,2.294) -- (3.800,2.592);
\node[cBar,anchor=east] at (3.752,2.802) {$\mathcal H'(s)$};
\node[gray,anchor=west] at (3.372,0.792) {$\frac r2s^2+us^4$};
\node[cWalk,anchor=west] at (0.000,3.312) {點線：$\frac{r'}2s^2+u's^4$};
\end{tikzpicture}}

\newcommand{\TwelveFigPairs}{%
\begin{tikzpicture}[font=\scriptsize,baseline=(current bounding box.center)]
\fill (0.000,0.000) circle (1.6pt);
\draw[gray] (-0.170,-0.200) -- (0.000,0.000) -- (0.170,-0.200);
\draw[cWalk,thick] (0.000,0.160) circle (0.160);
\node at (0,-0.55) {type 1};
\fill (1.300,0.000) circle (1.6pt);
\draw[gray] (1.130,-0.200) -- (1.300,0.000) -- (1.470,-0.200);
\fill (2.100,0.000) circle (1.6pt);
\draw[gray] (1.930,-0.200) -- (2.100,0.000) -- (2.270,-0.200);
\draw[cWalk,thick] (1.3,0) to[bend left=45] (2.1,0);
\draw[cWalk,thick] (1.3,0) to[bend right=45] (2.1,0);
\node at (1.7,-0.55) {type 2};
\fill (3.200,0.000) circle (1.6pt);
\draw[gray] (3.030,-0.200) -- (3.200,0.000) -- (3.370,-0.200);
\fill (4.000,0.000) circle (1.6pt);
\draw[gray] (3.830,-0.200) -- (4.000,0.000) -- (4.170,-0.200);
\fill (3.600,0.600) circle (1.6pt);
\draw[gray] (3.430,0.400) -- (3.600,0.600) -- (3.770,0.400);
\draw[cWalk,thick] (3.2,0) -- (4.0,0) -- (3.6,0.6) -- cycle;
\node at (3.6,-0.55) {type 3};
\draw[gray!40,rounded corners=2pt] (4.880,0.230) rectangle (6.160,0.950);
\fill (5.000,0.550) circle (1.6pt);
\fill (5.260,0.550) circle (1.6pt);
\fill (5.520,0.550) circle (1.6pt);
\draw[cWalk,thick] (5.520,0.460) circle (0.090);
\fill (5.780,0.550) circle (1.6pt);
\draw[cWalk,thick] (5.780,0.460) circle (0.090);
\fill (6.040,0.550) circle (1.6pt);
\draw[cWalk,thick] (6.040,0.460) circle (0.090);
\draw[cWalk,thick] (5.000,0.550) to[bend left=50] (5.260,0.550);
\draw[cWalk,thick] (5.000,0.550) to[bend left=20] (5.260,0.550);
\draw[gray!40,rounded corners=2pt] (6.300,0.230) rectangle (7.580,0.950);
\fill (6.420,0.550) circle (1.6pt);
\fill (6.680,0.550) circle (1.6pt);
\draw[cWalk,thick] (6.680,0.460) circle (0.090);
\fill (6.940,0.550) circle (1.6pt);
\fill (7.200,0.550) circle (1.6pt);
\draw[cWalk,thick] (7.200,0.460) circle (0.090);
\fill (7.460,0.550) circle (1.6pt);
\draw[cWalk,thick] (7.460,0.460) circle (0.090);
\draw[cWalk,thick] (6.420,0.550) to[bend left=50] (6.940,0.550);
\draw[cWalk,thick] (6.420,0.550) to[bend left=20] (6.940,0.550);
\draw[gray!40,rounded corners=2pt] (7.720,0.230) rectangle (9.000,0.950);
\fill (7.840,0.550) circle (1.6pt);
\fill (8.100,0.550) circle (1.6pt);
\draw[cWalk,thick] (8.100,0.460) circle (0.090);
\fill (8.360,0.550) circle (1.6pt);
\draw[cWalk,thick] (8.360,0.460) circle (0.090);
\fill (8.620,0.550) circle (1.6pt);
\fill (8.880,0.550) circle (1.6pt);
\draw[cWalk,thick] (8.880,0.460) circle (0.090);
\draw[cWalk,thick] (7.840,0.550) to[bend left=50] (8.620,0.550);
\draw[cWalk,thick] (7.840,0.550) to[bend left=20] (8.620,0.550);
\draw[gray!40,rounded corners=2pt] (9.140,0.230) rectangle (10.420,0.950);
\fill (9.260,0.550) circle (1.6pt);
\fill (9.520,0.550) circle (1.6pt);
\draw[cWalk,thick] (9.520,0.460) circle (0.090);
\fill (9.780,0.550) circle (1.6pt);
\draw[cWalk,thick] (9.780,0.460) circle (0.090);
\fill (10.040,0.550) circle (1.6pt);
\draw[cWalk,thick] (10.040,0.460) circle (0.090);
\fill (10.300,0.550) circle (1.6pt);
\draw[cWalk,thick] (9.260,0.550) to[bend left=50] (10.300,0.550);
\draw[cWalk,thick] (9.260,0.550) to[bend left=20] (10.300,0.550);
\draw[gray!40,rounded corners=2pt] (10.560,0.230) rectangle (11.840,0.950);
\fill (10.680,0.550) circle (1.6pt);
\draw[cWalk,thick] (10.680,0.460) circle (0.090);
\fill (10.940,0.550) circle (1.6pt);
\fill (11.200,0.550) circle (1.6pt);
\fill (11.460,0.550) circle (1.6pt);
\draw[cWalk,thick] (11.460,0.460) circle (0.090);
\fill (11.720,0.550) circle (1.6pt);
\draw[cWalk,thick] (11.720,0.460) circle (0.090);
\draw[cWalk,thick] (10.940,0.550) to[bend left=50] (11.200,0.550);
\draw[cWalk,thick] (10.940,0.550) to[bend left=20] (11.200,0.550);
\draw[gray!40,rounded corners=2pt] (4.880,-0.770) rectangle (6.160,-0.050);
\fill (5.000,-0.450) circle (1.6pt);
\draw[cWalk,thick] (5.000,-0.540) circle (0.090);
\fill (5.260,-0.450) circle (1.6pt);
\fill (5.520,-0.450) circle (1.6pt);
\draw[cWalk,thick] (5.520,-0.540) circle (0.090);
\fill (5.780,-0.450) circle (1.6pt);
\fill (6.040,-0.450) circle (1.6pt);
\draw[cWalk,thick] (6.040,-0.540) circle (0.090);
\draw[cWalk,thick] (5.260,-0.450) to[bend left=50] (5.780,-0.450);
\draw[cWalk,thick] (5.260,-0.450) to[bend left=20] (5.780,-0.450);
\draw[gray!40,rounded corners=2pt] (6.300,-0.770) rectangle (7.580,-0.050);
\fill (6.420,-0.450) circle (1.6pt);
\draw[cWalk,thick] (6.420,-0.540) circle (0.090);
\fill (6.680,-0.450) circle (1.6pt);
\fill (6.940,-0.450) circle (1.6pt);
\draw[cWalk,thick] (6.940,-0.540) circle (0.090);
\fill (7.200,-0.450) circle (1.6pt);
\draw[cWalk,thick] (7.200,-0.540) circle (0.090);
\fill (7.460,-0.450) circle (1.6pt);
\draw[cWalk,thick] (6.680,-0.450) to[bend left=50] (7.460,-0.450);
\draw[cWalk,thick] (6.680,-0.450) to[bend left=20] (7.460,-0.450);
\draw[gray!40,rounded corners=2pt] (7.720,-0.770) rectangle (9.000,-0.050);
\fill (7.840,-0.450) circle (1.6pt);
\draw[cWalk,thick] (7.840,-0.540) circle (0.090);
\fill (8.100,-0.450) circle (1.6pt);
\draw[cWalk,thick] (8.100,-0.540) circle (0.090);
\fill (8.360,-0.450) circle (1.6pt);
\fill (8.620,-0.450) circle (1.6pt);
\fill (8.880,-0.450) circle (1.6pt);
\draw[cWalk,thick] (8.880,-0.540) circle (0.090);
\draw[cWalk,thick] (8.360,-0.450) to[bend left=50] (8.620,-0.450);
\draw[cWalk,thick] (8.360,-0.450) to[bend left=20] (8.620,-0.450);
\draw[gray!40,rounded corners=2pt] (9.140,-0.770) rectangle (10.420,-0.050);
\fill (9.260,-0.450) circle (1.6pt);
\draw[cWalk,thick] (9.260,-0.540) circle (0.090);
\fill (9.520,-0.450) circle (1.6pt);
\draw[cWalk,thick] (9.520,-0.540) circle (0.090);
\fill (9.780,-0.450) circle (1.6pt);
\fill (10.040,-0.450) circle (1.6pt);
\draw[cWalk,thick] (10.040,-0.540) circle (0.090);
\fill (10.300,-0.450) circle (1.6pt);
\draw[cWalk,thick] (9.780,-0.450) to[bend left=50] (10.300,-0.450);
\draw[cWalk,thick] (9.780,-0.450) to[bend left=20] (10.300,-0.450);
\draw[gray!40,rounded corners=2pt] (10.560,-0.770) rectangle (11.840,-0.050);
\fill (10.680,-0.450) circle (1.6pt);
\draw[cWalk,thick] (10.680,-0.540) circle (0.090);
\fill (10.940,-0.450) circle (1.6pt);
\draw[cWalk,thick] (10.940,-0.540) circle (0.090);
\fill (11.200,-0.450) circle (1.6pt);
\draw[cWalk,thick] (11.200,-0.540) circle (0.090);
\fill (11.460,-0.450) circle (1.6pt);
\fill (11.720,-0.450) circle (1.6pt);
\draw[cWalk,thick] (11.460,-0.450) to[bend left=50] (11.720,-0.450);
\draw[cWalk,thick] (11.460,-0.450) to[bend left=20] (11.720,-0.450);
\node at (8.440,-1.050) {Ex.~12.9：$\frac{5!}{3!\,2!}=10$ 張片段圖};
\end{tikzpicture}}

\newcommand{\TwelveTabCensus}{%
\begin{tabular}{@{}lrr@{\quad}lrr@{}}
\toprule
\multicolumn{3}{c}{$n=3$} & \multicolumn{3}{c}{$n=4$}\\
環的大小 & 配對 & 圖 & 環的大小 & 配對 & 圖\\
\midrule
$(1,1,1)$ & 1 & 1 & $(1,1,1,1)$ & 1 & 1\\
$(2,1)$ & 6 & 3 & $(2,1,1)$ & 12 & 6\\
$(3)$ & 8 & 1 & $(3,1)$ & 32 & 4\\
 & &  & $(2,2)$ & 12 & 3\\
 & &  & $(4)$ & 48 & 3\\
\midrule
共 & 15 & & 共 & 105 & \\
\bottomrule
\end{tabular}}

\newcommand{\TwelveTabSfour}{%
\begin{tabular}{@{}lrrr@{}}
\toprule
$(n_1,n_2)$ & $n$ & $\frac{n!}{n_1!\,n_2!\,2^{n_2}}$ & 列舉\\
\midrule
$(2,1)$ & 4 & 6 & 6\,$\checkmark$\\
$(0,2)$ & 4 & 3 & 3\,$\checkmark$\\
$(3,1)$ & 5 & 10 & 10\,$\checkmark$\\
$(4,1)$ & 6 & 15 & 15\,$\checkmark$\\
$(2,2)$ & 6 & 45 & 45\,$\checkmark$\\
$(0,3)$ & 6 & 15 & 15\,$\checkmark$\\
\bottomrule
\end{tabular}}

\newcommand{\TwelveFigFlow}{%
\begin{tikzpicture}[font=\scriptsize,>=Stealth,baseline=(current bounding box.center)]
\draw[->,gray] (0,0) -- (7.090,0) node[right,black] {$r$};
\draw[->,gray] (3.420,0) -- (3.420,4.390) node[above,black] {$u$};
\draw[gray] (0.380,0) -- ++(0,-0.06) node[below,black] {$-0.4$};
\draw[gray] (1.900,0) -- ++(0,-0.06) node[below,black] {$-0.2$};
\draw[gray] (3.420,0) -- ++(0,-0.06) node[below,black] {$0$};
\draw[gray] (4.940,0) -- ++(0,-0.06) node[below,black] {$0.2$};
\draw[gray] (6.460,0) -- ++(0,-0.06) node[below,black] {$0.4$};
\draw[gray] (3.420,0.920) -- ++(-0.06,0) node[left,black] {$0.01$};
\draw[gray] (3.420,1.840) -- ++(-0.06,0) node[left,black] {$0.02$};
\draw[gray] (3.420,2.760) -- ++(-0.06,0) node[left,black] {$0.03$};
\draw[gray] (3.420,3.680) -- ++(-0.06,0) node[left,black] {$0.04$};
\draw[cBar,->,postaction={decorate},decoration={markings,mark=at position 0.5 with {\arrow{Stealth[length=3pt]}}}] (3.107,0.184) -- (3.089,0.191) -- (3.070,0.198) -- (3.050,0.205) -- (3.029,0.212) -- (3.007,0.220) -- (2.982,0.228) -- (2.957,0.236) -- (2.930,0.244) -- (2.900,0.253) -- (2.869,0.262) -- (2.836,0.271) -- (2.801,0.280) -- (2.763,0.290) -- (2.723,0.300) -- (2.680,0.310) -- (2.634,0.321) -- (2.584,0.331) -- (2.532,0.342) -- (2.476,0.353) -- (2.415,0.365) -- (2.351,0.376) -- (2.282,0.388) -- (2.209,0.400) -- (2.130,0.413) -- (2.045,0.425) -- (1.955,0.438) -- (1.858,0.450) -- (1.755,0.463) -- (1.644,0.476) -- (1.525,0.488) -- (1.398,0.501) -- (1.262,0.513) -- (1.116,0.525) -- (0.960,0.537) -- (0.793,0.548) -- (0.615,0.558) -- (0.423,0.567) -- (0.219,0.575) -- (-0.000,0.582);
\fill[cBar] (3.107,0.184) circle (0.9pt);
\draw[cExact,thick,postaction={decorate},decoration={markings,mark=at position 0.5 with {\arrow{Stealth[length=3pt]}}}] (3.259,0.184) -- (3.220,0.233) -- (3.175,0.292) -- (3.122,0.365) -- (3.063,0.451) -- (2.997,0.551) -- (2.926,0.665) -- (2.852,0.789) -- (2.778,0.922) -- (2.706,1.057) -- (2.639,1.189) -- (2.580,1.313) -- (2.528,1.424) -- (2.485,1.520) -- (2.449,1.600) -- (2.421,1.666) -- (2.400,1.717) -- (2.383,1.757) -- (2.370,1.788) -- (2.361,1.811) -- (2.354,1.829) -- (2.349,1.842) -- (2.345,1.851) -- (2.342,1.858) -- (2.340,1.863) -- (2.338,1.867) -- (2.337,1.870) -- (2.336,1.872) -- (2.336,1.874) -- (2.335,1.875) -- (2.335,1.875) -- (2.335,1.876) -- (2.335,1.876) -- (2.335,1.877) -- (2.335,1.877) -- (2.335,1.877) -- (2.335,1.877) -- (2.335,1.877) -- (2.335,1.877) -- (2.335,1.878);
\fill[cExact] (3.259,0.184) circle (0.9pt);
\draw[cBar,->,postaction={decorate},decoration={markings,mark=at position 0.5 with {\arrow{Stealth[length=3pt]}}}] (3.411,0.184) -- (3.418,0.191) -- (3.426,0.199) -- (3.436,0.207) -- (3.447,0.215) -- (3.459,0.224) -- (3.472,0.233) -- (3.487,0.242) -- (3.504,0.252) -- (3.523,0.262) -- (3.544,0.272) -- (3.567,0.282) -- (3.592,0.293) -- (3.620,0.305) -- (3.652,0.316) -- (3.686,0.329) -- (3.724,0.341) -- (3.766,0.354) -- (3.812,0.367) -- (3.862,0.381) -- (3.918,0.396) -- (3.979,0.411) -- (4.046,0.426) -- (4.120,0.442) -- (4.200,0.458) -- (4.288,0.475) -- (4.385,0.493) -- (4.490,0.511) -- (4.606,0.530) -- (4.732,0.549) -- (4.870,0.569) -- (5.020,0.590) -- (5.185,0.612) -- (5.364,0.634) -- (5.560,0.657) -- (5.774,0.681) -- (6.007,0.706) -- (6.261,0.733) -- (6.538,0.760) -- (6.840,0.788);
\fill[cBar] (3.411,0.184) circle (0.9pt);
\draw[cBar,->,postaction={decorate},decoration={markings,mark=at position 0.5 with {\arrow{Stealth[length=3pt]}}}] (1.598,3.680) -- (1.638,3.491) -- (1.671,3.329) -- (1.697,3.190) -- (1.719,3.068) -- (1.736,2.960) -- (1.749,2.864) -- (1.759,2.779) -- (1.765,2.701) -- (1.767,2.631) -- (1.767,2.566) -- (1.763,2.507) -- (1.757,2.452) -- (1.747,2.402) -- (1.735,2.354) -- (1.719,2.309) -- (1.701,2.267) -- (1.679,2.227) -- (1.654,2.189) -- (1.626,2.153) -- (1.595,2.117) -- (1.559,2.083) -- (1.520,2.050) -- (1.477,2.018) -- (1.430,1.986) -- (1.379,1.954) -- (1.323,1.923) -- (1.262,1.891) -- (1.195,1.860) -- (1.123,1.828) -- (1.045,1.795) -- (0.961,1.762) -- (0.870,1.728) -- (0.771,1.693) -- (0.665,1.657) -- (0.551,1.619) -- (0.428,1.580) -- (0.296,1.538) -- (0.153,1.495) -- (-0.000,1.448);
\fill[cBar] (1.598,3.680) circle (0.9pt);
\draw[cExact,thick,postaction={decorate},decoration={markings,mark=at position 0.5 with {\arrow{Stealth[length=3pt]}}}] (1.750,3.680) -- (1.986,2.864) -- (2.109,2.490) -- (2.182,2.281) -- (2.228,2.153) -- (2.260,2.069) -- (2.281,2.013) -- (2.296,1.974) -- (2.307,1.947) -- (2.314,1.927) -- (2.320,1.914) -- (2.324,1.904) -- (2.327,1.897) -- (2.329,1.891) -- (2.330,1.888) -- (2.331,1.885) -- (2.332,1.883) -- (2.333,1.881) -- (2.333,1.880) -- (2.333,1.880) -- (2.334,1.879) -- (2.334,1.879) -- (2.334,1.878) -- (2.334,1.878) -- (2.334,1.878) -- (2.334,1.878) -- (2.334,1.878) -- (2.334,1.878) -- (2.334,1.878) -- (2.334,1.878) -- (2.334,1.878) -- (2.334,1.878) -- (2.334,1.878) -- (2.334,1.878) -- (2.335,1.878) -- (2.335,1.878) -- (2.335,1.878) -- (2.335,1.878) -- (2.336,1.878) -- (2.337,1.878);
\fill[cExact] (1.750,3.680) circle (0.9pt);
\draw[cBar,->,postaction={decorate},decoration={markings,mark=at position 0.5 with {\arrow{Stealth[length=3pt]}}}] (1.902,3.680) -- (1.970,3.486) -- (2.031,3.325) -- (2.087,3.190) -- (2.140,3.076) -- (2.190,2.977) -- (2.239,2.892) -- (2.287,2.819) -- (2.334,2.754) -- (2.381,2.698) -- (2.429,2.649) -- (2.478,2.607) -- (2.528,2.569) -- (2.581,2.537) -- (2.635,2.509) -- (2.693,2.485) -- (2.754,2.465) -- (2.818,2.449) -- (2.887,2.435) -- (2.961,2.425) -- (3.040,2.418) -- (3.126,2.415) -- (3.218,2.413) -- (3.317,2.415) -- (3.425,2.420) -- (3.542,2.427) -- (3.668,2.438) -- (3.806,2.451) -- (3.956,2.468) -- (4.119,2.488) -- (4.297,2.511) -- (4.491,2.537) -- (4.702,2.567) -- (4.933,2.600) -- (5.184,2.637) -- (5.460,2.679) -- (5.760,2.724) -- (6.088,2.774) -- (6.447,2.829) -- (6.840,2.889);
\fill[cBar] (1.902,3.680) circle (0.9pt);
\draw[fill=white,thick] (3.420,0.000) circle (2.2pt);
\draw[fill=white,thick] (2.334,1.878) circle (2.2pt);
\draw[gray] (2.274,1.938) -- (0.684,4.232) node[above,black] {WF $\big(-\frac{1}{7},\frac{1}{49}\big)$};
\node[anchor=north west] at (3.470,-0.250) {Gaussian $(0,0)$};
\end{tikzpicture}}

\newcommand{\TwelveTabFixed}{%
\begin{tabular}{@{}lllrr@{}}
\toprule
$\epsilon$ & $r^*=-\frac\epsilon{6+\epsilon}$ & $u^*=\frac\epsilon{(6+\epsilon)^2}$ & 一階 $-\frac\epsilon6$ & $\frac\epsilon{36}$\\
\midrule
$1$ & $-\frac{1}{7}=-0.1429$ & $\frac{1}{49}=0.0204$ & $-0.1667$ & $0.0278$\\
$\frac{1}{2}$ & $-\frac{1}{13}=-0.0769$ & $\frac{2}{169}=0.0118$ & $-0.0833$ & $0.0139$\\
$\frac{1}{10}$ & $-\frac{1}{61}=-0.0164$ & $\frac{10}{3721}=0.0027$ & $-0.0167$ & $0.0028$\\
\bottomrule
\end{tabular}}

\newcommand{\TwelveFigU}{%
\begin{tikzpicture}[font=\scriptsize,>=Stealth,baseline=(current bounding box.center)]
\draw[->,gray] (0,0) -- (5.170,0) node[right,black] {$\ell$};
\draw[->,gray] (0.000,0) -- (0.000,3.970) node[above,black] {$u(\ell)$};
\draw[gray] (0.000,0) -- ++(0,-0.06) node[below,black] {$0$};
\draw[gray] (0.820,0) -- ++(0,-0.06) node[below,black] {$1$};
\draw[gray] (1.640,0) -- ++(0,-0.06) node[below,black] {$2$};
\draw[gray] (2.460,0) -- ++(0,-0.06) node[below,black] {$3$};
\draw[gray] (3.280,0) -- ++(0,-0.06) node[below,black] {$4$};
\draw[gray] (4.100,0) -- ++(0,-0.06) node[below,black] {$5$};
\draw[gray] (4.920,0) -- ++(0,-0.06) node[below,black] {$6$};
\draw[gray] (0.000,0.620) -- ++(-0.06,0) node[left,black] {$0.01$};
\draw[gray] (0.000,1.240) -- ++(-0.06,0) node[left,black] {$0.02$};
\draw[gray] (0.000,1.860) -- ++(-0.06,0) node[left,black] {$0.03$};
\draw[gray] (0.000,2.480) -- ++(-0.06,0) node[left,black] {$0.04$};
\draw[gray] (0.000,3.100) -- ++(-0.06,0) node[left,black] {$0.05$};
\draw[gray] (0.000,3.720) -- ++(-0.06,0) node[left,black] {$0.06$};
\draw[cBar,thick] (0.000,0.310) -- (0.126,0.351) -- (0.252,0.396) -- (0.378,0.445) -- (0.505,0.497) -- (0.631,0.554) -- (0.757,0.613) -- (0.883,0.675) -- (1.009,0.739) -- (1.135,0.804) -- (1.262,0.871) -- (1.388,0.937) -- (1.514,1.002) -- (1.640,1.065) -- (1.766,1.127) -- (1.892,1.185) -- (2.018,1.240) -- (2.145,1.292) -- (2.271,1.340) -- (2.397,1.383) -- (2.523,1.423) -- (2.649,1.459) -- (2.775,1.492) -- (2.902,1.521) -- (3.028,1.547) -- (3.154,1.569) -- (3.280,1.590) -- (3.406,1.607) -- (3.532,1.623) -- (3.658,1.636) -- (3.785,1.648) -- (3.911,1.658) -- (4.037,1.667) -- (4.163,1.675) -- (4.289,1.681) -- (4.415,1.687) -- (4.542,1.692) -- (4.668,1.696) -- (4.794,1.700) -- (4.920,1.703);
\fill[cBar] (0.820,0.644) circle (1.3pt);
\fill[cBar] (1.640,1.065) circle (1.3pt);
\fill[cBar] (4.100,1.671) circle (1.3pt);
\node[cBar,anchor=west] at (0.12,0.410) {$u_0=0.005$};
\draw[cWalk,thick] (0.000,3.720) -- (0.126,3.192) -- (0.252,2.846) -- (0.378,2.604) -- (0.505,2.426) -- (0.631,2.293) -- (0.757,2.189) -- (0.883,2.108) -- (1.009,2.043) -- (1.135,1.990) -- (1.262,1.947) -- (1.388,1.911) -- (1.514,1.882) -- (1.640,1.857) -- (1.766,1.837) -- (1.892,1.819) -- (2.018,1.805) -- (2.145,1.793) -- (2.271,1.782) -- (2.397,1.773) -- (2.523,1.766) -- (2.649,1.760) -- (2.775,1.754) -- (2.902,1.750) -- (3.028,1.746) -- (3.154,1.742) -- (3.280,1.739) -- (3.406,1.737) -- (3.532,1.735) -- (3.658,1.733) -- (3.785,1.731) -- (3.911,1.730) -- (4.037,1.729) -- (4.163,1.728) -- (4.289,1.727) -- (4.415,1.726) -- (4.542,1.726) -- (4.668,1.725) -- (4.794,1.725) -- (4.920,1.725);
\fill[cWalk] (0.820,2.146) circle (1.3pt);
\fill[cWalk] (1.640,1.857) circle (1.3pt);
\fill[cWalk] (4.100,1.728) circle (1.3pt);
\node[cWalk,anchor=west] at (0.12,3.840) {$u_0=0.06$};
\draw[gray,densely dashed] (0,1.722) -- (4.920,1.722) node[right,black] {$u^*=\frac1{36}$};
\draw[cExact,thick] (4.100,1.571) -- ++(0.1,0.1) -- ++(-0.1,0.1) -- ++(-0.1,-0.1) -- cycle;
\node[cExact,anchor=north] at (4.100,1.551) {$g_0(e^{-5})$};
\end{tikzpicture}}

\newcommand{\TwelveFigNu}{%
\begin{tikzpicture}[font=\scriptsize,>=Stealth,baseline=(current bounding box.center)]
\draw[->,gray] (0,0) -- (5.450,0) node[right,black] {$\epsilon$};
\draw[->,gray] (0.000,0) -- (0.000,2.850) node[above,black] {$\nu$};
\draw[gray] (0.000,0) -- ++(0,-0.06) node[below,black] {$0$};
\draw[gray] (2.600,0) -- ++(0,-0.06) node[below,black] {$0.5$};
\draw[gray] (5.200,0) -- ++(0,-0.06) node[below,black] {$1$};
\draw[gray] (0.000,0.000) -- ++(-0.06,0) node[left,black] {$0.50$};
\draw[gray] (0.000,1.300) -- ++(-0.06,0) node[left,black] {$0.55$};
\draw[gray] (0.000,2.600) -- ++(-0.06,0) node[left,black] {$0.60$};
\draw[cBar,thick] (0.000,0.000) -- (0.179,0.075) -- (0.359,0.149) -- (0.538,0.224) -- (0.717,0.299) -- (0.897,0.374) -- (1.076,0.448) -- (1.255,0.523) -- (1.434,0.598) -- (1.614,0.672) -- (1.793,0.747) -- (1.972,0.822) -- (2.152,0.897) -- (2.331,0.971) -- (2.510,1.046) -- (2.690,1.121) -- (2.869,1.195) -- (3.048,1.270) -- (3.228,1.345) -- (3.407,1.420) -- (3.586,1.494) -- (3.766,1.569) -- (3.945,1.644) -- (4.124,1.718) -- (4.303,1.793) -- (4.483,1.868) -- (4.662,1.943) -- (4.841,2.017) -- (5.021,2.092) -- (5.200,2.167);
\draw[cWalk,thick,densely dashed] (0.000,0.000) -- (0.179,0.075) -- (0.359,0.151) -- (0.538,0.228) -- (0.717,0.306) -- (0.897,0.385) -- (1.076,0.464) -- (1.255,0.545) -- (1.434,0.627) -- (1.614,0.709) -- (1.793,0.793) -- (1.972,0.877) -- (2.152,0.963) -- (2.331,1.050) -- (2.510,1.137) -- (2.690,1.226) -- (2.869,1.316) -- (3.048,1.408) -- (3.228,1.500) -- (3.407,1.594) -- (3.586,1.688) -- (3.766,1.784) -- (3.945,1.882) -- (4.124,1.980) -- (4.303,2.080) -- (4.483,2.181) -- (4.662,2.284) -- (4.841,2.388) -- (5.021,2.493) -- (5.200,2.600);
\draw[cExact,thick,densely dotted] (0.000,0.000) -- (0.179,0.073) -- (0.359,0.141) -- (0.538,0.205) -- (0.717,0.265) -- (0.897,0.321) -- (1.076,0.372) -- (1.255,0.420) -- (1.434,0.464) -- (1.614,0.505) -- (1.793,0.541) -- (1.972,0.574) -- (2.152,0.604) -- (2.331,0.631) -- (2.510,0.654) -- (2.690,0.675) -- (2.869,0.692) -- (3.048,0.707) -- (3.228,0.719) -- (3.407,0.728) -- (3.586,0.736) -- (3.766,0.740) -- (3.945,0.743) -- (4.124,0.743) -- (4.303,0.742) -- (4.483,0.738) -- (4.662,0.733) -- (4.841,0.726) -- (5.021,0.717) -- (5.200,0.707);
\node[cWalk,anchor=west] at (5.250,2.600) {不展開 $\frac1{2-\epsilon/3}$};
\node[cBar,anchor=west] at (5.250,2.117) {一階 $\frac12+\frac\epsilon{12}$};
\node[cExact,anchor=west] at (5.250,0.587) {矩陣 $1/\lambda_r$};
\end{tikzpicture}}

\newcommand{\TwelveTabNu}{%
\begin{tabular}{@{}lrrr@{}}
\toprule
$\epsilon$ & 一階 & 不展開 & 矩陣\\
\midrule
$\frac{1}{10}$ & 0.5083 & 0.5085 & 0.5076\\
$\frac{1}{2}$ & 0.5417 & 0.5455 & 0.5256\\
$1$ & 0.5833 & 0.6000 & 0.5272\\
\bottomrule
\end{tabular}}

\newcommand{\TwelveTabOn}{%
\begin{tabular}{@{}crrrlrrr@{}}
\toprule
$n$ & $2n+4$ & $A=4(n+2)$ & $B=4(n+8)$ & 一階 (12.68) & 不展開 & 矩陣 & 已知（引用）\\
\midrule
0 & 4 & 8 & 32 & $\frac{9}{16}=0.5625$ & 0.5714 & 0.5205 & 0.588\\
1 & 6 & 12 & 36 & $\frac{7}{12}=0.5833$ & 0.6000 & 0.5272 & 0.630\\
2 & 8 & 16 & 40 & $\frac{3}{5}=0.6000$ & 0.6250 & 0.5325 & 0.672\\
3 & 10 & 20 & 44 & $\frac{27}{44}=0.6136$ & 0.6471 & 0.5368 & 0.711\\
\bottomrule
\end{tabular}}

\newcommand{\TwelveTabPc}{%
\begin{tabular}{@{}lrrrrr@{}}
\toprule
$d$ & $d-3$ & 右邊 & 數值 & 書上 $2d-3$ & 書上右邊\\
\midrule
$2$ & $-1$ & $-1$ & $-1.0000$ & {\color{cWrong}$1$} & {\color{cWrong}$-3$}\\
$\frac{5}{2}$ & $-\frac{1}{2}$ & $-\frac{1}{2}$ & $-0.5000$ & {\color{cWrong}$2$} & {\color{cWrong}$-\frac{5}{2}$}\\
$3$ & $0$ & $0$ & -- & {\color{cWrong}$3$} & {\color{cWrong}$-2$}\\
$4$ & $1$ & $1$ & -- & {\color{cWrong}$5$} & {\color{cWrong}$-1$}\\
\bottomrule
\end{tabular}}

\newcommand{\TwelveTabErrata}{%
\begin{tabular}{@{}lll@{}}
\toprule
位置 & 書上 & 應為\\
\midrule
p.~293，(12.17) 前 & {\color{cWrong}Figure 12.16(b)} & {\color{cExact}Figure 12.11(b)}\\
p.~295，(12.32) & {\color{cWrong}$N\to N'e^{d\delta\ell}\approx N'(1+\delta\ell)$} & {\color{cExact}$\approx N'(1+d\,\delta\ell)$}\\
p.~295，(12.33) & {\color{cWrong}$\vec q\to\vec q\,'e^{-d\delta\ell}$} & {\color{cExact}$\vec q\to\vec q\,'e^{-\delta\ell}$}\\
pp.~296--297，(12.36)(12.38)(12.39) & {\color{cWrong}$(4-d)u-36u^2$} & {\color{cExact}$(4-d)u-36d\,u^2$}\\
p.~298，(12.47) & {\color{cWrong}$\int_0^\ell u\,d\ell'=12d\int\cdots$} & {\color{cExact}$12d\int_0^\ell u\,d\ell'=12d\int\cdots$}\\
p.~298，(12.48) & {\color{cWrong}$12d\int_{u_0}^u\frac{du'}{\epsilon u'-36du'^2}$} & {\color{cExact}$12d\int_{u_0}^u\frac{u'\,du'}{\epsilon u'-36du'^2}$}\\
p.~299，Ex.~12.4 & {\color{cWrong}(12.53) follows from (12.51)} & {\color{cExact}(12.52)}\\
p.~299，(12.60) 後 & {\color{cWrong}3D Ising $\nu\sim0.61$} & {\color{cExact}約 0.630（文獻值）}\\
p.~315，(12.89) 後 & {\color{cWrong}$u^2r^{2d-3}=r^{-1}(ur^{-(4-d)/2})^2$} & {\color{cExact}$u^2r^{d-3}=r\,(ur^{-(4-d)/2})^2$}\\
p.~320 & {\color{cWrong}$5!/(2!\,1!)$（第二次）} & {\color{cExact}$5!/(3!\,2!)=10$}\\
p.~322，(12:S-7) & {\color{cWrong}$\mathcal T_1^{n_1}\mathcal T_3^{n_3}\mathcal T_3^{n_3}$，$n_1\ge1$} & {\color{cExact}$\mathcal T_1^{n_1}\mathcal T_2^{n_2}\mathcal T_3^{n_3}$，$n_1\ge0$}\\
\bottomrule
\end{tabular}}

\newcommand{\TwelveTabBook}{%
\begin{tabular}{@{}rllp{10.5cm}@{}}
\toprule
章 & 主角 & 動畫 & 關鍵結果與數字\\
\midrule
1 & 一維漫步 & W101--W112 & \raggedright $\langle x^2\rangle=N$；$P_{10}(0)=\frac{252}{1024}=0.2461$，高斯近似 0.2523\tabularnewline
2 & 生成函數與格子 & W201--W212 & \raggedright 二維回原點 $\binom N{N/2}^2$：1、4、36、400、4900、63504\tabularnewline
3 & 回家問題 & W301--W308 & \raggedright Pólya：一、二維一定回家；簡單立方晶格的回家機率 0.3405\tabularnewline
4 & 牆與邊界 & W401--W408 & \raggedright 吸收牆與鏡像法；$x_0=2$、$L=5$ 先到左牆的機率 $\frac{3}{5}$\tabularnewline
5 & 偏向與持續 & W501--W507 & \raggedright 持續長度 $l_p=\frac Q{1-Q}l$；$Q=\frac12$、$N=1000$：$\langle x^2\rangle=2996$\tabularnewline
6 & 漫步的形狀 & W601--W607 & \raggedright 非球度 $A_d=\frac{4+2d}{4+5d}$，$A_3=\frac{10}{19}=0.5263$\tabularnewline
7 & 自迴避 & W701--W706 & \raggedright $c_{14}=2374444$；Flory $\nu_F=\frac3{d+2}$\tabularnewline
8 & 長尾與標度 & W801--W804 & \raggedright 上臨界維度 $d_c=4$；二維的連結常數 $\mu=2.638$\tabularnewline
9 & 臨界現象 & W901--W906 & \raggedright O($n$)，$n\to0$＝自迴避；Flory 0.600 對 0.588（差 2\%，0.0204）\tabularnewline
10 & 受限的高分子 & W1001--W1005 & \raggedright $\rho=|r|/2u$；$\nu'=\frac{1}{2}$\tabularnewline
11 & 重整化群 & W1101--W1106 & \raggedright Koch 曲線 $\frac{\ln4}{\ln3}=1.2619$；$\nu=1/\Lambda_r$\tabularnewline
12 & $\epsilon$ 展開 & W1201--W1206 & \raggedright $\nu=\frac12+\frac{(n+2)\epsilon}{4(n+8)}$：Ising 0.5833、自迴避 0.5625\tabularnewline
\midrule
 & 全書 & 87 幕 & 從 $\sqrt N$ 到 $N^{0.588}$（一階 $N^{9/16}$）\\
\bottomrule
\end{tabular}}


\section{\texorpdfstring{重整化群（續）與 $\epsilon$ 展開（第 12 章）}{重整化群（續）與 ε 展開（第 12 章）}}

\begin{keybox}[先看這裡]
第 11 章的動量殼只有二次項（Gaussian 模型）：各個 Fourier 模態互不相干，積掉外圈殼層只給一個常數，所以只得到平均場的 $\nu=\frac12$。
這一章加上四次項 $u\,S^4$，照同樣四步走一遍：
(1)~把布里淵區外圈一層薄薄的\textbf{殼層}裡的模態先積分掉（\textbf{殼層積分}）；
(2)~四次項把殼層和內區的模態纏在一起，對殼層做 Gaussian 平均＝把殼層的「手」兩兩牽起來，而取對數之後\textbf{只剩連通的圖}（\textbf{連結叢集}定理）；
(3)~重新標度回原來的樣子，得到 $r$、$u$ 兩個參數的微分方程——參數的\textbf{流}；
(4)~四維以下（$\epsilon=4-d>0$）流有一個新的\textbf{不動點}（Wilson--Fisher 點），它離 Gaussian 點只有 $O(\epsilon)$，所以可以把 $\epsilon$ 當小參數展開臨界指數：
Ising $\nu=\frac12+\frac\epsilon{12}$，O($n$) $\nu=\frac12+\frac{(n+2)\epsilon}{4(n+8)}$；$n=0$ 就是自迴避漫步，$\epsilon=1$ 得 $\nu\approx\frac9{16}$（已知 0.588）。
最後兩節（12.5--12.6）用費曼圖講同一件事，只取概念。這是全書最後一章，本節最後附一頁\textbf{全書回顧}。

\textbf{記號}：$\epsilon=4-d$；GLW 的 $r$（約化溫度）、$u$（四次耦合）；布里淵區取半徑 $\Lambda=e^{-\ell}$ 的球；殼層裡 $r+q^2$ 換成 $r+1$。
\textbf{本系列把書上 (12.37) 的 $d$ 吸收進 $u$}（本系列的 $u$ ＝書上的 $ud$；書上 (12.66)(12.67) 本來就這樣寫，$\nu$ 不受影響）。
\begin{center}
\small
\begin{tabular}{@{}lll@{\qquad}lll@{}}
\toprule
節 & 內容 & 動畫 & 節 & 內容 & 動畫\\
\midrule
12.1--12.2.4 & 四次項、殼層 & W1201 & 12.2.5--6 & $\epsilon$ 展開求 $\nu$ & W1204\\
12.2.4、12.7 & 配對、連結叢集 & W1202 & 12.3--12.4 & O($n$)、自迴避漫步 & W1205\\
12.2.4--5 & 一迴圈、$(r,u)$ 的流 & W1203 & 12.5--12.6 & 圖解法、全書回顧 & W1206\\
\bottomrule
\end{tabular}
\end{center}
數字一律由 \texttt{codes/utils/walk\_eps.py} 重算，逐式校對見 \texttt{ref/md/ch12\_formulas.md}，
圖表由 \texttt{ch12/make\_figs.py} 產生並先做數字檢查。連續播放：\texttt{Ch12Demo}（W1201--W1206）。
\end{keybox}

% =====================================================================
\subsection{加上四次項：殼層與「怎樣正確地積掉自由度」}
\leavevmode\bk{pp.~285--290}\ \anim{W1201}\par\nopagebreak

\textbf{問}：上一章只有二次項；加上四次項之後，還能「積掉一半」嗎？

\textbf{模型}（(12.1)--(12.4)）：一分量（Ising 型）的 GLW Hamiltonian 寫在 Fourier 空間，
\[\mathcal H[S(\vec q)]=\sum_{\vec q}\frac12\big(r+q^2\big)S(\vec q)S(-\vec q)+\frac uN\sum_{\vec q_1\ldots\vec q_4}S(\vec q_1)S(\vec q_2)S(\vec q_3)S(\vec q_4)\,\delta_{\vec q_1+\vec q_2+\vec q_3+\vec q_4}.\tag{12.3--12.4}\]
二次項裡每個 $\vec q$ 各管各的；四次項的 $\delta$ 是動量守恆：四個 $\vec q$ 加起來是 0 才有貢獻，於是不同的模態被綁在一起。
\textbf{倒晶格}（Fig.~12.1--12.2）：格距 $a$ 的晶格 $\leftrightarrow$ 格距 $2\pi/a$ 的倒晶格；系統邊長 $L\propto N^{1/d}$ 越大，允許的 $\vec q$ 越密（間距 $2\pi/L$）。$N$ 個格點對應 $N$ 個 $\vec q$。

\textbf{殼層}（Fig.~12.3--12.4，圖 \ref{twelve-fig-shell} 左）：臨界點附近只有小的 $q$（長波長）重要，所以把立方的布里淵區換成半徑 $\Lambda=e^{-\ell}$ 的球；
每一步積掉 $e^{-\ell-\delta\ell}<|\vec q|<e^{-\ell}$ 這一層。\textbf{殼層積分}的白話：每一個 $S(\vec q)$ 都是一個積分變數，$Z$ 是對全部 $S(\vec q)$ 的積分；
先把殼層裡那幾個變數積完，剩下的就是內區變數的有效 Hamiltonian。殼層占的模態比例是體積比
\[\frac{\sum_{\rm shell}1}{N}=1-e^{-d\,\delta\ell}\approx d\,\delta\ell,\qquad dq=e^{-\ell}-e^{-\ell-\delta\ell}\approx e^{-\ell}\delta\ell .\tag{12.20--12.24}\]
\textbf{算到底}：$d=3$、$\delta\ell=0.1$：0.2592 對 0.3000；$\delta\ell=0.01$：0.0296 對 0.0300；$d=4$、$\delta\ell=0.01$：0.0392 對 0.0400——$\delta\ell$ 越小越準。
真的在半徑 200 的球裡數格點也一樣（0.2591、0.0296，差不到 $2\times10^{-4}$）。$\ell=0.5$（$\Lambda=0.6065$）、$\delta\ell=0.01$：$dq=0.006035$，近似 $e^{-\ell}\delta\ell=0.006065$。

\textbf{怎樣正確地積掉自由度}（(12.5)--(12.10)）：把變數分成留下的 $S_i$ 與要積掉的 $S_i'$：
\[\begin{gathered}Z(K)=\sum_{\{S_i\},\{S_i'\}}e^{-H(K,S_i,S_i')},\qquad A(K')\,e^{-H'(K',S_i)}=\sum_{\{S_i'\}}e^{-H(K,S_i,S_i')},\\
Z_N(K)=A(K')\,Z_{N'}(K'),\qquad K'=\mathcal R(K).\end{gathered}\tag{12.5--12.10}\]
要點：$H'$ 必須和 $H$ \textbf{同一個形式}，只是參數 $K\to K'$；跑出來的常數 $A$ 只進自由能。
\textbf{兩個變數的玩具}（書外）：$s$ 在內區、$S'$ 在殼層，只留 (12.11) 的 $6u$ 頂點：
$\mathcal H=\frac r2s^2+us^4+\frac{r+1}2S'^2+6u\,s^2S'^2$。對 $S'$ 是 Gaussian 積分，$\int e^{-\alpha S'^2}dS'=\sqrt{\pi/\alpha}$，$\alpha=\frac{r+1}2+6us^2$：
\[A=\sqrt{\frac{2\pi}{r+1}},\qquad \mathcal H'(s)=\frac r2s^2+us^4+\frac12\ln\Big(1+\frac{12u\,s^2}{r+1}\Big).\]
用 $\ln(1+x)=x-\frac{x^2}2+\frac{x^3}3-\cdots$ 展開：$s^2$ 的係數多了 $\frac{6u}{1+r}$、$s^4$ 的係數少了 $\frac{36u^2}{(1+r)^2}$，還冒出一個 $s^6$ 項：
\[r'=r+\frac{12u}{1+r},\qquad u'=u-\frac{36u^2}{(1+r)^2},\qquad c_6=\frac16\Big(\frac{12u}{1+r}\Big)^3.\]
$r=0.2$、$u=0.01$（$\frac{12u}{1+r}=0.1$）：$r'=0.3000$、$u'=0.0075$、$c_6=0.000167$（書上丟掉的高階項）；$\mathcal H'(1)=0.1577$（閉式＝數值積分）。
配分函數：二重積分 $Z=9.0517$，而 $A=2.2882$、$Z'=\int e^{-\mathcal H'(s)}ds=3.9558$，$A\,Z'=9.0517$——正是 $Z_N=A(K')Z_{N'}(K')$。
\textbf{12 和 36 就是下一節要在殼層裡數出來的數字}（真正的殼層還多一個 $d\,\delta\ell$）。

\begin{figure}[htbp]
\centering
\begin{minipage}[c]{0.28\linewidth}\centering\TwelveFigShell\end{minipage}\hfill
\begin{minipage}[c]{0.31\linewidth}\centering\footnotesize\setlength{\tabcolsep}{3pt}\TwelveTabShell\end{minipage}\hfill
\begin{minipage}[c]{0.38\linewidth}\centering\TwelveFigToy\end{minipage}
\caption{左：球形布里淵區（平面示意），點是允許的 $\vec q$，殼層的模態（橙）先積掉。中：殼層比例（數格點：半徑 200 的球）。右：兩變數玩具（$r=0.2$、$u=0.01$），精確的 $\mathcal H'(s)$ 與截斷的 $\frac{r'}2s^2+u's^4$ 幾乎重合}\label{twelve-fig-shell}
\end{figure}

\textbf{四次項按幾隻腳在殼層分類}（(12.11)）：$\binom4k$ 給係數 $1,4,6,4,1$（0 到 4 隻腳在殼層）；殼層很薄時，三、四隻腳在殼層的兩項可略（書上的說法）。
\textbf{重新標度}（(12.29)--(12.34)）：積完殼層，布里淵區半徑變小；把 $\vec q\to\vec q\,'e^{-\delta\ell}$ 讓它變回原大小，$N\to N'e^{d\delta\ell}$（模態變少），再令 $S\to S''e^{\delta\ell}$ 讓 $q^2$ 項的係數維持 1：
$q^2S^2\to q'^2e^{-2\delta\ell}S''^2e^{2\delta\ell}$（乘 1）、$rS^2\to re^{2\delta\ell}S''^2$、$\frac uNS^4\to\frac{u\,e^{4\delta\ell}}{N'e^{d\delta\ell}}S''^4$。
\textbf{算到底}（$d=3$、$\delta\ell=0.1$）：$q^2$ 項 1.0000、$r$ 項 $e^{2\delta\ell}=1.2214$、$u$ 項 $e^{(4-d)\delta\ell}=1.1052$。$d<4$ 時 $u$ 被放大——四次項在長波長變重要，這就是平均場在四維以下失效的原因。
\begin{errbox}[p.~295，(12.32)(12.33)]
(12.32) 印成 $N\to N'e^{d\delta\ell}\approx N'(1+\delta\ell)$，應為 $N'(1+d\,\delta\ell)$：$d=3$、$\delta\ell=0.01$ 時 $e^{0.03}=1.0305$，不是 1.0100。
(12.33) 印成 $\vec q\to\vec q\,'e^{-d\delta\ell}\approx\vec q\,'(1-d\,\delta\ell)$，應為 $\vec q\,'e^{-\delta\ell}$（正文自己也這樣說）：照書上 $q^2$ 項要乘 $e^{-2(d-1)\delta\ell}=0.6703$（$d=3$、$\delta\ell=0.1$），模型就不再同形。
\end{errbox}

\textbf{揭曉}：積掉殼層、重新標度之後，模型還是 GLW 的樣子，只剩 $r$ 與 $u$ 兩個數在變。

% =====================================================================
\subsection{配對遊戲：只剩連通的圖}
\leavevmode\bk{pp.~289--292、317--322}\ \anim{W1202}\par\nopagebreak

\textbf{問}：殼層的變數積掉之後，四次頂點會變成什麼？

\textbf{直觀}：把四次項裡碰到殼層的部分當成微擾 $V$，$e^{-V}=\sum_n\frac{(-V)^n}{n!}$，再對殼層變數做 Gaussian 平均。Gaussian 平均有一條規則（第 6 章的 Wick 配對）：
\textbf{一串變數乘積的平均＝把它們兩兩配對的所有方法之和}，每一對給
\[\big\langle S(\vec Q_1)S(\vec Q_2)\big\rangle=\frac{\delta(\vec Q_1+\vec Q_2)}{r+Q_1^2}\qquad(\text{殼層 }Q\approx1\text{：}\tfrac1{r+1}).\tag{12:S-2}\]
把頂點畫成一個點加四隻手（Fig.~12.6）：一隻手在殼層的 (a) 係數 $\frac{4u}N$、兩隻手在殼層的 (b) 係數 $\frac{6u}N$（(12.14)(12.15)）。殼層的手兩兩牽起來，內區的手留著，變成新的 $r$、$u$。

\textbf{片段與環}（補充 12.7，圖 \ref{twelve-fig-pairs}）：每個 (b) 片段有兩隻殼層的手；全部牽好之後，片段連成一個個環：
type 1（自己的兩手相牽）、type 2（兩個片段雙線）、type 3（三角環）……$n$ 個片段共有 $(2n-1)!!$ 種牽法（3、15、105、945）。
\textbf{算到底}：$n=3$：$15=1+6+8$（三個 type 1／type 1＋type 2／三角環），不同的片段圖 1、3、1 張；三角環的 8 種牽法是 $2^3$。
\textbf{Ex.~12.9}：5 個片段配成 3 個 type 1 ＋ 1 個 type 2。945 種牽法裡有 20 種＝10 張片段圖 $\times$ 2 種雙線接法（圖 \ref{twelve-fig-pairs} 上右的 10 張）：
\[\frac{5!}{3!\,2!}=10,\qquad\text{一般}\quad\frac{n!}{n_1!\,n_2!\,(2!)^{n_2}}\quad(n=n_1+2n_2).\tag{12:S-4}\]
理由：$n$ 個片段排成一列有 $n!$ 種；type 1 之間互換（$n_1!$）、type 2 之間互換（$n_2!$）、每個 type 2 內部互換（$2!$）都是同一張圖。圖 \ref{twelve-fig-pairs} 右下的表對 $n\le6$ 全部和列舉一致。
\begin{errbox}[p.~320]
同一頁前面寫對了 $\frac{5!}{3!\,2!}$，第二次卻寫成 $\frac{5!}{2!\,1!}$——照書上是 60，列舉只有 10 張圖。
\end{errbox}

\textbf{為什麼只剩連通的圖}（(12:S-3)--(12:S-7)）：記 $\mathcal T_k$ 為「一個 type $k$ 環」的值。把 $\frac1{n!}$ 和上面的組合因子相乘，$n!$ 正好消掉：
\[\sum_{n_1,n_2}\frac{\mathcal T_1^{n_1}\mathcal T_2^{n_2}}{n_1!\,n_2!}=e^{\mathcal T_1+\mathcal T_2},\qquad
\sum_{n_1,n_2,n_3,\ldots}\frac{\mathcal T_1^{n_1}\mathcal T_2^{n_2}\mathcal T_3^{n_3}}{n_1!\,n_2!\,n_3!}\cdots=\exp\big[\mathcal T_1+\mathcal T_2+\mathcal T_3+\cdots\big].\tag{12:S-5--12:S-7}\]
\textbf{所有圖的和＝$\exp$（連通圖的和）}：取對數就只剩連通的圖——不連通的圖由指數自動算好。一般的因子是 $\frac{n!}{\prod_kn_k!\,f_k^{n_k}}$，
$f_k$＝換頂點而配對不變的方法數：1、2、6、8、10（$f_3=6=3!$，Ex.~12.11；〔注〕$k\ge4$ 時 $f_k=2k$，不是 $k!$）。
\begin{errbox}[p.~322，(12:S-7)]
分子印成 $\mathcal T_1^{n_1}\mathcal T_3^{n_3}\mathcal T_3^{n_3}$、$n_1$ 從 1 起加；應為 $\mathcal T_1^{n_1}\mathcal T_2^{n_2}\mathcal T_3^{n_3}$，$n_1$ 也從 0 起（$n_1=0$ 的項不能少，否則得不到 $e^{\mathcal T_1}$）。
\end{errbox}
\textbf{最簡形}：一個 Gaussian 變數 $\langle s^{2n}\rangle=(2n-1)!!$ 就是所有配對；連成一個環的配對有 $w_k=2^{k-1}(k-1)!$ 種（1、2、8、48）。於是
$\sum_n(2n-1)!!\frac{x^n}{n!}=\exp\big(\sum_kw_k\frac{x^k}{k!}\big)=(1-2x)^{-1/2}$；$x=0.1$：三種算法都是 1.1180。

\textbf{數值核對}（Ex.~12.1）：$X=-as^2$，$s$ 的變異數 $\sigma^2=1$，$a=0.1$。$\langle e^X\rangle$ 用四種算法：數值積分、所有圖的和 $\sum_n\frac{\langle X^n\rangle}{n!}$、
$\exp$（連通圖）$=\exp\sum_k\frac{\kappa_k}{k!}$、閉式 $(1+2a\sigma^2)^{-1/2}$，全部 0.9129（對數 $-0.0912$）。連通圖（環）$\kappa_k=(-a\sigma^2)^k2^{k-1}(k-1)!$：$-0.1000$、0.0200、$-0.0080$、0.0048。
二階：全部 $\frac{\langle X^2\rangle}2=\frac{3a^2}2=0.0150$ ＝ 連通 $\frac{\kappa_2}2=0.0100$ ＋ 不連通 $\frac{\kappa_1^2}2=0.0050$。三個殼層變數（耦合矩陣）也一樣：$\det(1+2A)^{-1/2}$＝三維求積＝環的和＝0.8087。

\begin{figure}[htbp]
\centering
\TwelveFigPairs\\[6pt]
\begin{minipage}[c]{0.55\linewidth}\centering\footnotesize\setlength{\tabcolsep}{3pt}\TwelveTabCensus\end{minipage}\hfill
\begin{minipage}[c]{0.43\linewidth}\centering\footnotesize\setlength{\tabcolsep}{4pt}\TwelveTabSfour\end{minipage}
\caption{配對。上左：三種環（黑點＝片段，灰＝內區的手，橙＝殼層的手牽成的線）。上右：Ex.~12.9 的 10 張片段圖（雙線那一對之外各自成圈）。
左下：$n=3,4$ 個片段的配對普查（配對數／不同的圖數）。右下：(12:S-4) 與列舉}\label{twelve-fig-pairs}
\end{figure}

\textbf{揭曉}：指數化把不連通的圖自動算好——只要畫連通的圖；Fig.~12.11 的兩張一迴圈圖就是下一節的主角。

% =====================================================================
\subsection{一迴圈與 \texorpdfstring{$(r,u)$}{(r,u)} 平面的流}
\leavevmode\bk{pp.~292--297}\ \anim{W1203}\par\nopagebreak

\textbf{問}：兩張一迴圈圖各給多少？

\textbf{(a) 改 $r$}（Fig.~12.11(a)，(12.16)--(12.18)）：一個 $6u$ 頂點，兩隻殼層的手自己接成一圈，剩兩隻內區的手——形狀就是二次項。
\textbf{(b) 改 $u$}（Fig.~12.11(b)，(12.19)）：二階，兩個頂點各用兩隻殼層的手互接（type 2），剩四隻內區的手——形狀就是四次項：
\[\text{(a)}=\frac{6u}{N(r+1)}\sum_{\vec Q\in\rm shell}\ \sum_{\vec q}S(\vec q)S(-\vec q),\qquad
\text{(b)}=\frac{36u^2}{N^2(r+1)^2}\sum_{\vec Q,\vec Q'\in\rm shell}\ \sum_{\vec q_1..\vec q_4}S\cdots S\,\delta .\tag{12.18--12.19}\]
係數的來歷正是 12.1 節的玩具：(a) 的 $6u$ 要和 $\frac r2$ 比，所以 $r$ 多 $2\times6=12$；(b) 是二階的 $\frac12\times6\times6\times2$（雙線兩種接法）$=36$，進四次項時帶負號（$\ln$ 的二階項）。
\textbf{殼層求和}：$\sum_{\vec Q\in\rm shell}1=N\frac{\text{殼層體積}}{\text{布里淵區體積}}=N\frac{d\,\delta q}\Lambda=N\,d\,\delta\ell$（(12.20)--(12.25)，就是 12.1 節的 0.0392 對 0.0400）；
(b) 的兩個殼層動量被動量守恆 $\vec Q+\vec Q'+\vec q_1+\vec q_2=0$ 綁住（(12.26)），取 $\vec q_1+\vec q_2\approx0$ 只剩一個殼層求和（(12.27)）。於是 (12.28)：
\[r\to r+\frac{12u}{r+1}d\,\delta\ell,\qquad u\to u-\frac{36u^2}{(r+1)^2}d\,\delta\ell .\]
\textbf{乘上重新標度的乘數}（$e^{2\delta\ell}$、$e^{(4-d)\delta\ell}$），令 $\delta\ell\to0$ 得遞迴方程（(12.35)--(12.38)；本系列把 $d$ 併進 $u$）：
\[\frac{dr}{d\ell}=2r+\frac{12u}{1+r},\qquad\frac{du}{d\ell}=\epsilon u-\frac{36u^2}{(1+r)^2}.\tag{12.37--12.38}\]
\begin{errbox}[pp.~296--297，(12.36)(12.38)(12.39)；p.~293]
書上這三式的 $u^2$ 項寫成 $36u^2$，少了 $d$：(12.28)、(12.35)、(12.40)、(12.41) 和 $u^*\to\epsilon/36d$ 都有 $d$，應為 $(4-d)u-36d\,u^2$。
（本系列把 $d$ 併進 $u$，兩種寫法的 $\nu$ 相同。）另外 p.~293 (12.17) 前的「Figure 12.16(b)」應為 Figure 12.11(b)。〔注〕(12.35) 第二個指數前少了負號。
\end{errbox}

\textbf{不動點}（$\frac{dr}{d\ell}=\frac{du}{d\ell}=0$）：Gaussian 點 $(0,0)$；另一個由 $\epsilon=\frac{36u}{(1+r)^2}$ 與 $2r(1+r)=-12u$ 解出 $\epsilon(1+r)=-6r$：
\[\text{Wilson--Fisher：}\quad r^*=-\frac{\epsilon}{6+\epsilon},\qquad u^*=\frac{\epsilon}{(6+\epsilon)^2}\qquad(\text{一階 }-\tfrac\epsilon6,\ \tfrac\epsilon{36}).\]
$\epsilon=1$：$(-\frac17,\frac1{49})=(-0.1429,0.0204)$（表見圖 \ref{twelve-fig-flow} 右）。$\epsilon\to0$ 時兩個不動點合而為一（書上 $u^*\to\epsilon/36d$）。
\textbf{Gaussian 點的穩定性}：在 $(0,0)$ 線性化，矩陣 $\begin{pmatrix}2&12\\0&\epsilon\end{pmatrix}$ 的特徵值是 $2$ 與 $\epsilon$（$\epsilon=1$：2.0000、1.0000）。
$\epsilon>0$ 時 $u$ 方向也是相關的——Gaussian 點不穩定，臨界的系統會被推向 WF 點；$d>4$（$\epsilon<0$）時 $u$ 無關，平均場正確——第 8 章的上臨界維度 4 從這裡來。

\textbf{流向圖}（圖 \ref{twelve-fig-flow} 左，$\epsilon=1$）：$u_0=0.002$ 與 0.04 各取三個起點 $r_0=r_c(u_0)$、$r_c\pm0.02$。中間那條（臨界線，綠）流進 WF 點；
旁邊兩條先靠近 WF 點，再沿相關方向往 $r\to+$（高溫相）或 $r\to-$（低溫相）跑掉（畫到 $|r|=0.45$）。不管 $u_0$ 是多少，臨界的系統最後都停在同一點——\textbf{普適性}。

\begin{figure}[htbp]
\centering
\begin{minipage}[c]{0.5\linewidth}\centering\TwelveFigFlow\end{minipage}\hfill
\begin{minipage}[c]{0.48\linewidth}\centering\footnotesize\setlength{\tabcolsep}{3pt}\TwelveTabFixed\end{minipage}
\caption{左：$(r,u)$ 平面的流（$\epsilon=1$；白圈是 Gaussian 點與 Wilson--Fisher 點）。右：WF 點的精確值與一階近似；$\epsilon\to0$ 時都趨向 $(0,0)$}\label{twelve-fig-flow}
\end{figure}

\textbf{揭曉}：四維以下，臨界的系統都流向 Wilson--Fisher 點——臨界指數由它附近的流決定。

% =====================================================================
\subsection{\texorpdfstring{$\epsilon$ 展開求 $\nu$}{ε 展開求 ν}}
\leavevmode\bk{pp.~296--300}\ \anim{W1204}\par\nopagebreak

\textbf{問}：耦合 $u$ 最後流到哪裡？$\nu$ 是多少？

\textbf{$u$ 的流}（(12.39)--(12.41)，Ex.~12.2）：在 $r=0$ 附近 $\frac{du}{d\ell}=\epsilon u-36u^2$——小的 $u$ 像 $e^{\epsilon\ell}$ 一樣長大、大的 $u$ 被 $-36u^2$ 拉下來。分離變數積分：
\[\int_{u_0}^{u(\ell)}\frac{du}{\epsilon u-36u^2}=\ell\quad\Rightarrow\quad u(\ell)=\frac{\epsilon u_0e^{\epsilon\ell}}{\epsilon+36u_0(e^{\epsilon\ell}-1)}\ \xrightarrow{\ \ell\to\infty\ }\ \frac{\epsilon}{36}.\tag{12.40--12.41}\]
\textbf{算到底}（$\epsilon=1$，圖 \ref{twelve-fig-u}）：$u_0=0.005$：$\ell=1,2,5$ 得 0.0104、0.0172、0.0270；$u_0=0.06$：0.0346、0.0300、0.0279——從兩邊都趨向 $u^*=\frac1{36}=0.0278$（閉式＝數值積分）。
\textbf{Ex.~12.3}（Fig.~12.13）：另一個不動點 $u=0$ 的斜率是 $+\epsilon$（不穩定）：$u_0=10^{-6}$ 時 $u(2)/u_0=7.3874\approx e^2=7.3891$；$u^*$ 的斜率是 $-\epsilon$（穩定）。

\begin{figure}[htbp]
\centering
\TwelveFigU
\caption{(12.41) 的 $u(\ell)$（$\epsilon=1$），點是數值積分；菱形是 12.6 節 (12.86) 的 $g_0(q_0=e^{-5})$，落在同一條曲線上}\label{twelve-fig-u}
\end{figure}

\textbf{$r$ 的流}（(12.42)--(12.51)）：把 $\frac{12u}{1+r}$ 展開，$\frac{dr}{d\ell}\approx2r+12u-12ur$。令 $r=t-6u$（書上 $t-6ud$，在保留 $d$ 的寫法下是對的）消掉常數項 $12u$，一階得
\[\frac{dt}{d\ell}=2t-12u\,t\quad\Rightarrow\quad t(\ell)=t(0)\,e^{2\ell-12\int_0^\ell u\,d\ell'}.\tag{12.45--12.46}\]
積分換成對 $u$：$d\ell'=\frac{du'}{\epsilon u'-36u'^2}$，所以 $12\int_0^\ell u\,d\ell'=12\int_{u_0}^{u}\frac{u'\,du'}{\epsilon u'-36u'^2}=-\frac13\ln(\epsilon-36u')\big|_{u_0}^{u}$，再代入 (12.41)：
\[t(\ell)=t(0)\,e^{2\ell}\Big[\frac{\epsilon}{\epsilon+36u_0(e^{\epsilon\ell}-1)}\Big]^{1/3}\ \xrightarrow{\ \ell\text{ 大}\ }\ \propto t(0)\,e^{(2-\epsilon/3)\ell}.\tag{12.49--12.51}\]
\textbf{算到底}：$t_0=10^{-4}$、$u_0=0.01$、$\ell=3$：閉式＝數值積分＝0.0203。$t$ 的指數增長率是 $2-\frac\epsilon3$，比 Gaussian 的 2 慢一點。
\begin{errbox}[p.~298，(12.47)(12.48)]
(12.47) 左邊 $\int_0^\ell u\,d\ell'$ 少了 $12d$。(12.48) 印成 $12d\int_{u_0}^u\frac{du'}{\epsilon u'-36du'^2}$，分子少了 $u'$，應為 $12d\int_{u_0}^u\frac{u'\,du'}{\epsilon u'-36du'^2}$：
$u_0=0.001\to u=0.005$（$\epsilon=1$、$d=3$）正確的是 0.2207，等於 (12.49)；照書上是 81.7804。
\end{errbox}

\textbf{讀出 $\nu$}（(12.52)--(12.59)）：$t(0)$ 是「離臨界的距離」，$t(0)\propto T-T_c^*$（$T_c^*$ 是被漲落移動過的臨界溫度）。一直流到 $r(\ell^*)=D$（一個普通大小的數，那時關聯長度約 1 個格距）就停：
\[e^{(2-\epsilon/3)\ell^*}=\frac D{t(0)}\ \Rightarrow\ e^{\ell^*}=\Big(\frac D{T-T_c^*}\Big)^{\frac1{2-\epsilon/3}},\qquad 1=\xi(\ell^*)=\xi(0)e^{-\ell^*}\ \Rightarrow\ \xi(0)\propto(T-T_c^*)^{-\frac1{2-\epsilon/3}}.\tag{12.52--12.58}\]
\[\nu=\frac1{2-\epsilon/3}\approx\frac12\Big(1+\frac\epsilon6\Big)=\frac12+\frac\epsilon{12},\qquad\epsilon=1:\ \nu=\frac7{12}=0.5833 .\tag{12.59--12.60}\]
平均場 0.5 $\to$ 一階 0.5833 $\to$ 三維 Ising 已知約 0.630（文獻值，只引用）：第一次「算」出不是 $\frac12$ 的臨界指數。
\begin{errbox}[p.~299，Ex.~12.4；(12.60) 後]
Ex.~12.4「Show that (12.53) follows from (12.51)」應為 (12.52)——(12.53) 在練習之後才定義。(12.60) 後說三維 Ising 的「$\nu\sim0.61$」，已知值約 0.630。
\end{errbox}
\textbf{臨界點被移動}（Ex.~12.4）：$t(0)=r(0)+6u_0$，所以臨界點在 $r_c\approx-6u_0$（書上 $-6u_0d$）。精確地說，臨界線是 Gaussian 點 $\epsilon$ 方向的特徵向量 $r_c\approx-\frac{12u_0}{2-\epsilon}$：
$u_0=10^{-4}$、$\epsilon=0.1$：二分法 $-0.000631$、特徵向量 $-0.000632$、一階 $-0.000600$（$u_0=10^{-3}$、$\epsilon=1$：$-0.0111$ 對 $-0.0120$）。

\textbf{$\nu$ 的三種讀法}（本系列，圖 \ref{twelve-fig-nu}）：一階 $\frac12+\frac\epsilon{12}$；(12.54) 不展開 $\frac1{2-\epsilon/3}$；把完整的一迴圈方程在 WF 點線性化、$\nu=1/\lambda_r$
（$\epsilon=1$ 的矩陣精確是 $\begin{pmatrix}5/3&14\\1/21&-1\end{pmatrix}$，特徵值 $\frac{1\pm\sqrt{22}}3=1.8968,\ -1.2301$，$\nu=0.5272$）。
三條曲線在 $\epsilon=0$ 相切，只有 $\epsilon$ 的一次項相同——一迴圈只準到 $O(\epsilon)$，$O(\epsilon^2)$ 要二迴圈的圖；$\epsilon=1$ 時三者相差不小，這是 12.5 節「隱憂」的第一條。

\begin{figure}[htbp]
\centering
\begin{minipage}[c]{0.62\linewidth}\centering\TwelveFigNu\end{minipage}\hfill
\begin{minipage}[c]{0.36\linewidth}\centering\footnotesize\setlength{\tabcolsep}{4pt}\TwelveTabNu\end{minipage}
\caption{$\nu$ 的三種讀法（Ising）：一階、不展開、線性化矩陣；在 $\epsilon=0$ 相切}\label{twelve-fig-nu}
\end{figure}

\textbf{揭曉}：Wilson--Fisher 點把 $\nu$ 從平均場的 $\frac12$ 推到 $\frac7{12}$；$\epsilon$ 是唯一的小參數。

% =====================================================================
\subsection{O(\texorpdfstring{$n$}{n}) 模型與自迴避漫步}
\leavevmode\bk{pp.~300--304}\ \anim{W1205}\par\nopagebreak

\textbf{問}：自旋有 $n$ 個分量時，12 和 36 會變成什麼？

\textbf{模型}（(12.61)、Fig.~12.14）：四次項 $\frac uN\sum(\vec S_1\cdot\vec S_2)(\vec S_3\cdot\vec S_4)\delta$，殼層的平均（Ex.~12.5）
\[\big\langle S^\alpha(\vec Q_1)S^\beta(\vec Q_2)\big\rangle=\frac{\delta_{\alpha\beta}}{1+r}\,\delta(\vec Q_1+\vec Q_2)\tag{12.63}\]
——不同分量不配對（殼層的權重 $e^{-\frac12(1+r)\sum_\alpha|S^\alpha|^2}$ 對各分量獨立）。〔注〕(12.63) 印成 $S^\beta(S(\vec Q_2)$，多一個「$S($」。
\textbf{數分量}（Fig.~12.15，Ex.~12.6）：四隻手選兩隻到殼層（6 種）：兩隻在\textbf{同一個}點積（2 種）$\to$ 繞一圈 $\sum_\beta\delta_{\beta\beta}=n$；在\textbf{不同}點積（4 種）$\to$ 1。合計
\[\text{(a)}+\text{(b)}=\frac{2n+4}{1+r}\,\frac uN\sum_{\vec Q}\sum_{\vec q}\vec S(\vec q)\cdot\vec S(-\vec q)\tag{12.64}\]
（$n=1$ 回到 6；〔注〕書上少了 $\frac uN$）。所以 $12\to A=2(2n+4)=4(n+2)$；同樣數 $u$ 的圖（Fig.~12.16--12.17，Ex.~12.7）得 $36\to B=4(n+8)$：
\[\frac{dr}{d\ell}=2r+\frac{4(n+2)u}{1+r},\qquad\frac{du}{d\ell}=\epsilon u-\frac{4(n+8)}{(1+r)^2}u^2 .\tag{12.66--12.67}\]
\textbf{讀出 $\nu$}：照 12.4 節，$u^*\approx\frac\epsilon B$、$r=t-\frac A2u$，$\frac{dt}{d\ell}=2t-Au\,t$，$t\propto e^{(2-A\epsilon/B)\ell}$：
\[\nu=\frac1{2-\frac{(n+2)\epsilon}{n+8}}\approx\frac12+\frac{(n+2)\epsilon}{4(n+8)},\qquad n=0:\ \nu=\frac12+\frac\epsilon{16},\quad\epsilon=1:\ \frac9{16}=0.5625 .\tag{12.68--12.70}\]
$n=1$ 回到 $\frac12+\frac\epsilon{12}$（$\frac AB=\frac13$）。〔注〕(12.70) 的 0.563 是 0.5625 四捨五入。$n=0$ 的 WF 點是 $(-\frac19,\frac2{81})$。
程式真的對 $n=1..4$ 做了張量縮併：$2n+4$ 與 $4(n+8)$ 都對（表 \ref{twelve-tab-on}）。

\begin{table}[htbp]
\centering\footnotesize\setlength{\tabcolsep}{4pt}
\TwelveTabOn
\caption{O($n$) 模型（$\epsilon=1$）：殼層的因子與 $\nu$ 的三種讀法；已知值只引用。$n=0$ 是自迴避漫步、$n=1$ Ising、$n=2$ XY、$n=3$ Heisenberg}\label{twelve-tab-on}
\end{table}

\textbf{Flory 與 $\epsilon$ 展開}（(12.71)(12.72)）：第 7 章的 Flory 指數 $\nu_F=\frac3{d+2}$（$d=3$：$\frac35=0.6$）寫成 $\epsilon$ 的函數再展開：
\[\nu_F=\frac3{6-\epsilon}=\frac12+\frac\epsilon{12}+\frac{\epsilon^2}{72}+\frac{\epsilon^3}{432}+\cdots\qquad(\text{係數 }\tfrac12(\tfrac16)^k).\tag{12.71--12.72}\]
一次係數 $\frac1{12}$ 對 $n=0$ 的 $\frac1{16}$：\textbf{Flory 不是精確的}（雖然數值接近：0.5833 對 0.5625）。〔注〕Ising 不展開的 $\frac1{2-\epsilon/3}$ 恰好就是 $\frac3{6-\epsilon}$——巧合。
\textbf{算到底}：$N=1000$ 步的漫步大小 $N^\nu$：普通 $\sqrt N=31.62$、一階 $N^{9/16}=48.70$、已知 $N^{0.588}=58.08$——自迴避讓鏈撐開將近一倍。

\textbf{$\epsilon$ 展開的四個隱憂}（p.~304）：$\epsilon=1$ 並不小；級數是漸近級數（收斂半徑 0）；可能有 $e^{-A/\epsilon}$ 這種冪級數抓不到的項；重整化群本身沒有被嚴格證明。
儘管如此，高階的 $\epsilon$ 展開（配合重新求和）給出的指數和實驗、數值模擬相符到小數第二、三位。

\textbf{揭曉}：同一套遞迴，$n\to0$ 就是自迴避漫步——第 1 章的 $\sqrt N$ 在這裡變成 $N^{0.5625}$（已知 $N^{0.588}$）。

% =====================================================================
\subsection{圖解法的重整化群（概念）}
\leavevmode\bk{pp.~304--317}\ \anim{W1206}\par\nopagebreak

\textbf{問}：不切殼層、直接用費曼圖，也能做重整化嗎？

\textbf{重整化傳播子}（(12.73)，Fig.~12.18--12.20）：所有「插一塊 $\Sigma$」的鏈加起來是幾何級數（和第 7 章 W704 的 Dyson 方程同一個形狀）：
\[\frac1{r+q^2}+\frac1{r+q^2}\,\Sigma\,\frac1{r+q^2}+\cdots=\frac1{r+q^2-\Sigma(r,\vec q)}.\tag{12.73}\]
$\Sigma$ 叫\textbf{質量算子}，整條鏈畫成一條雙線。把圖裡內部的動量依大小排序，只算 $q<q_0<q_1,q_2,q_3$ 的部分（Fig.~12.21 的「西瓜」圖），
截斷 $q_0$ 就扮演殼層半徑的角色；重整化傳播子寫成 $q^2+\Sigma'=q_0^{2-\eta}\mathcal S(q/q_0)$，$y=1/\nu$（(12.75)--(12.80)）。

\textbf{四點函數與 $g_0$}（(12.81)--(12.86)）：把有效的四次耦合寫成 $u_R=q_0^{4-d+2\eta}g$，定義 $g_0(q_0)=q_0^{-4+d-2\eta}u_R$，得
\[q_0\frac{dg_0}{dq_0}=(-4+d-2\eta)g_0+\mathcal Bg_0^2 .\tag{12.86}\]
換變數 $q_0=e^{-\ell}$（截斷往小動量走＝$\ell$ 變大），$q_0\frac d{dq_0}=-\frac d{d\ell}$：$\frac{dg_0}{d\ell}=(\epsilon+2\eta)g_0-\mathcal Bg_0^2$——$\eta=0$、$\mathcal B=36$ 就是 (12.39)。
\textbf{算到底}：$g_0(1)=0.005$ 積到 $q_0=e^{-5}=0.0067$：0.0270，正是 12.4 節的 $u(\ell=5)$（圖 \ref{twelve-fig-u} 的菱形）；不動點 $g^*=\frac\epsilon{\mathcal B}=0.0278$。

\textbf{冪次數算}（(12.87)--(12.89)，Fig.~12.27）：直接做微擾，第 $n$ 階含 $(u\,r^{-(4-d)/2})^n$，$d<4$、$r\to0$ 時每一階都發散——普通的微擾在臨界點失效。
但重整化的耦合 $u_R\propto r^{(4-d)/2}$ 剛好把它抵消，有效的展開參數停在 $g^*$。以西瓜圖為例：特徵動量 $Q\propto r^{1/2}$，三個 $d^dQ$ 給 $r^{3d/2}$、$\delta$ 給 $r^{-d/2}$、三個 $G=\frac1{r+Q^2}$ 給 $r^{-3}$：
\[\int u^2G(Q_1)G(Q_2)G(Q_3)\,d^dQ_1d^dQ_2d^dQ_3\,\delta(\vec Q_1+\vec Q_2+\vec Q_3-\vec q)\ \propto\ u^2r^{d-3}=r\,\big(ur^{-(4-d)/2}\big)^2\tag{12.89}\]
（$\Sigma$ 和 $r$ 同單位，所以前面是 $r$）。數值積分量出的冪次與 $d-3$ 一致：
\begin{center}\footnotesize\setlength{\tabcolsep}{5pt}\TwelveTabPc\end{center}
\begin{errbox}[p.~315，(12.89) 後]
書上寫 $u^2r^{2d-3}=r^{-1}\times(ur^{-(4-d)/2})^2$，兩邊本來就不相等（$2d-3$ 對 $d-5$）。應為 $u^2r^{d-3}=r\times(ur^{-(4-d)/2})^2$：$d=2$ 數值冪次 $-1.0000$，照書上是 $+1$。
\end{errbox}
\textbf{反項與 Callan--Symanzik 方程}（Fig.~12.28，pp.~315--317，只用文字）：截斷 $q_0$ 是算法的產物，物理的答案不能依賴它；要求「改變 $q_0$ 而答案不變」就得到一條對 $q_0$ 的微分方程——又是重整化群方程。

\textbf{揭曉}：切殼層（Wilson）與費曼圖（場論）走不同的路，得到同一條流方程與同一個不動點。

% =====================================================================
\subsection{練習 12.1--12.11 解答}
{\small
\begin{enumerate}[label=\textbf{12.\arabic*},itemsep=0.05em]
\item 連結叢集。\quad 答：$n$ 個片段的 Gaussian 平均＝所有配對的和；每個配對把片段分成環（type $k$），同一種分法的數目 $\frac{n!}{\prod_kn_k!\,f_k^{n_k}}$，
配上 $e^{-V}$ 展開的 $\frac1{n!}$ 就是 $\prod_k\frac{\mathcal T_k^{n_k}}{n_k!}$，對所有 $n_k$ 求和 $\Rightarrow\exp(\sum_k\mathcal T_k)$。數值：$X=-as^2$ 時 0.9129（四種算法）、二階 $0.0150=0.0100+0.0050$；三個變數 0.8087。
\item 解 (12.39)。\quad 答：部分分式 $\frac1{\epsilon u-36du^2}=\frac1\epsilon\big(\frac1u+\frac{36d}{\epsilon-36du}\big)$ $\Rightarrow\frac1\epsilon\ln\frac{u(\epsilon-36du_0)}{u_0(\epsilon-36du)}=\ell$，解出 $u$ 就是 (12.41)；
$\ell\to\infty$ 得 $\frac\epsilon{36d}$。數值（本系列 $d$ 併進 $u$）$u_0=0.005$：0.0104、0.0172、0.0270。
\item 另一個不動點。\quad 答：$u=0$ 附近 $\frac{d\,\delta u}{d\ell}=\epsilon\,\delta u$ $\Rightarrow\delta u\propto e^{\epsilon\ell}$ 變大（不穩定）；$u^*$ 處斜率 $\epsilon-72du^*=-\epsilon$（穩定）$\Rightarrow$ Fig.~12.13 的箭頭都指向 $u^*$。
$u_0=10^{-6}$：$u(2)/u_0=7.3874\approx e^2$。
\item 由 (12.51) 得 (12.52)（題目的 (12.53) 應為 (12.52)）。\quad 答：大 $\ell$ 時 $\big[\frac\epsilon{\epsilon+36du_0(e^{\epsilon\ell}-1)}\big]^{1/3}\to\big(\frac\epsilon{36du_0}\big)^{1/3}e^{-\epsilon\ell/3}$ $\Rightarrow t\propto t(0)e^{(2-\epsilon/3)\ell}$；
$r^*(0)=t(0)=r(0)+6u_0d\propto T-T_c^*$ $\Rightarrow$ 臨界點移到 $r_c\approx-6u_0d$（本系列 $-6u_0$；數值 $-0.000631$、特徵向量 $-0.000632$、一階 $-0.000600$）。
\item (12.63)。\quad 答：殼層的權重 $\exp[-\frac12(1+r)\sum_\alpha S^\alpha(\vec Q)S^\alpha(-\vec Q)]$ 對各分量獨立，不同分量的平均是 0，同分量就是一分量的 (12:S-2) $\Rightarrow\frac{\delta_{\alpha\beta}}{1+r}\delta(\vec Q_1+\vec Q_2)$。
\item (12.64)。\quad 答：四隻手選兩隻在殼層（6 種）：同一個點積（2 種）$\sum_\beta\delta_{\beta\beta}=n$，不同點積（4 種）各 1 $\Rightarrow2n+4$；和 $\frac r2$ 比較多一個 2 $\Rightarrow A=4(n+2)$（$n=1$：12）。
\item (12.65)。\quad 答：兩個頂點各選兩隻殼層的手、互接（$6\times6\times2=72$ 種）；對外面四隻手對稱化後是 $(\vec S\cdot\vec S)^2$ 的 $8(n+8)$ 倍（$n=1$：72），乘二階的 $\frac12$ $\Rightarrow4(n+8)$（$n=1$：36；einsum 核對 $n=1..4$）。
\item (12.68)。\quad 答：把 12 換成 $4(n+2)$、36 換成 $4(n+8)$：$t\propto e^{(2-\frac{(n+2)\epsilon}{n+8})\ell}$ $\Rightarrow\nu=\frac1{2-(n+2)\epsilon/(n+8)}\approx\frac12+\frac{(n+2)\epsilon}{4(n+8)}$；$n=0$、$\epsilon=1$：$\frac9{16}=0.5625$。
\item $\frac{5!}{3!\,2!}$。\quad 答：先排 5 個片段（$5!$），前三個各自成 type 1、後兩個成 type 2；前三個互換（$3!$）、後兩個互換（$2!$）不產生新的圖 $\Rightarrow10$。
列舉：945 種配對裡 20 種＝10 張圖 $\times$ 2 種雙線接法；書上第二次寫的 $\frac{5!}{2!\,1!}=60$ 是筆誤。
\item (12:S-4)。\quad 答：$n=n_1+2n_2$ 個片段排成一列：$n!$；type 1 之間互換 $n_1!$、type 2 之間互換 $n_2!$、每個 type 2 內部互換 $2!$ $\Rightarrow\frac{n!}{n_1!\,n_2!\,(2!)^{n_2}}$（$n\le6$ 全部＝列舉，圖 \ref{twelve-fig-pairs}）。
\item $f_3=6$。\quad 答：三角環的 $3!=6$ 種換頂點都保持「每個頂點和另外兩個各連一線」$\Rightarrow f_3=6$（三角環的腳配對 8 種＝$2^3$）。〔注〕$k\ge4$ 的環 $f_k=2k$（只有轉與翻），例如 $f_4=8$。
\end{enumerate}}

% =====================================================================
\subsection{第 12 章的回顧與筆誤一覽}
\begin{center}
\footnotesize\setlength{\tabcolsep}{4pt}
\begin{tabular}{@{}lll@{}}
\toprule
主題（書上的式子） & 結果 & 動畫\\
\midrule
殼層與積掉自由度 (12.1)--(12.34) & 玩具 $r'=r+\frac{12u}{1+r}=0.3000$、$u'=u-\frac{36u^2}{(1+r)^2}=0.0075$；$u$ 乘 $e^{\epsilon\delta\ell}$ & W1201\\
配對、連結叢集 (12:S-1)--(12:S-7) & 只剩連通圖；$\frac{5!}{3!\,2!}=10$、$f_3=6$；$\langle e^X\rangle=0.9129$ & W1202\\
一迴圈與流 (12.16)--(12.38) & $\frac{dr}{d\ell}=2r+\frac{12u}{1+r}$、$\frac{du}{d\ell}=\epsilon u-\frac{36u^2}{(1+r)^2}$；WF 點 $(-\frac17,\frac1{49})$ & W1203\\
$\epsilon$ 展開 (12.39)--(12.60) & $u\to\frac\epsilon{36}$；$\nu=\frac12+\frac\epsilon{12}=0.5833$（已知 0.630） & W1204\\
O($n$) 與自迴避 (12.61)--(12.72) & $\nu=\frac12+\frac{(n+2)\epsilon}{4(n+8)}$；$n=0$：0.5625（已知 0.588）；Flory 不精確 & W1205\\
圖解法 (12.73)--(12.89) & (12.86)＝(12.39)（$q_0=e^{-\ell}$）；$u^2r^{d-3}=r(ur^{-(4-d)/2})^2$ & W1206\\
\bottomrule
\end{tabular}
\end{center}
\textbf{筆誤一覽}（給 2004 年版；〔注〕見各節）：下表在文中都用紅框標出，每一條都用程式核對過。
\begin{center}
\footnotesize\renewcommand{\arraystretch}{0.9}
\TwelveTabErrata
\end{center}

% =====================================================================
\clearpage
\subsection{全書回顧}

十二章、87 幕動畫走的是同一條路：從一枚硬幣出發（第 1 章，$\langle x^2\rangle=N$、距離只長成 $\sqrt N$），用\textbf{生成函數}把數路線變成做積分（第 2--6 章：格子、回家、牆、偏向、形狀）；
再加上「不能走回頭路」，漫步變成真正的難題（第 7--8 章：自迴避、Flory 的 $\frac3{d+2}$、四維是分界）；
接著發現自迴避漫步就是 $n\to0$ 的磁鐵（第 9--10 章），漫步的指數就是臨界指數；最後用\textbf{重整化群}——放大、積掉、再放大——
在四維以下找到 Wilson--Fisher 不動點，把 $\nu$ 從 $\frac12$ 一路算到 $\frac12+\frac\epsilon{16}$（第 11--12 章）。下表每章一列，數字與動畫 W1206 的全書回顧相同。
\begin{center}
\small\setlength{\tabcolsep}{4pt}\renewcommand{\arraystretch}{1.3}
\TwelveTabBook
\end{center}
\textbf{一句話}：漫步、磁鐵與重整化群說的是同一個故事——自迴避的鏈在三維裡 $R\sim N^{0.588}$，一迴圈的 $\epsilon$ 展開給 $N^{0.5625}$，Flory 給 $N^{0.6}$，而沒有限制的漫步永遠是 $\sqrt N$。
